% Options for packages loaded elsewhere
% Options for packages loaded elsewhere
\PassOptionsToPackage{unicode}{hyperref}
\PassOptionsToPackage{hyphens}{url}
\PassOptionsToPackage{dvipsnames,svgnames,x11names}{xcolor}
%
\documentclass[
  english,
  11pt,
  a4paper,
  authoryear,
  preprint]{elsarticle}
\usepackage{xcolor}
\usepackage[margin=25mm]{geometry}
\usepackage{amsmath,amssymb}
\setcounter{secnumdepth}{5}
\usepackage{iftex}
\ifPDFTeX
  \usepackage[T1]{fontenc}
  \usepackage[utf8]{inputenc}
  \usepackage{textcomp} % provide euro and other symbols
\else % if luatex or xetex
  \usepackage{unicode-math} % this also loads fontspec
  \defaultfontfeatures{Scale=MatchLowercase}
  \defaultfontfeatures[\rmfamily]{Ligatures=TeX,Scale=1}
\fi
\usepackage{lmodern}
\ifPDFTeX\else
  % xetex/luatex font selection
\fi
% Use upquote if available, for straight quotes in verbatim environments
\IfFileExists{upquote.sty}{\usepackage{upquote}}{}
\IfFileExists{microtype.sty}{% use microtype if available
  \usepackage[]{microtype}
  \UseMicrotypeSet[protrusion]{basicmath} % disable protrusion for tt fonts
}{}
\makeatletter
\@ifundefined{KOMAClassName}{% if non-KOMA class
  \IfFileExists{parskip.sty}{%
    \usepackage{parskip}
  }{% else
    \setlength{\parindent}{0pt}
    \setlength{\parskip}{6pt plus 2pt minus 1pt}}
}{% if KOMA class
  \KOMAoptions{parskip=half}}
\makeatother
% Make \paragraph and \subparagraph free-standing
\makeatletter
\ifx\paragraph\undefined\else
  \let\oldparagraph\paragraph
  \renewcommand{\paragraph}{
    \@ifstar
      \xxxParagraphStar
      \xxxParagraphNoStar
  }
  \newcommand{\xxxParagraphStar}[1]{\oldparagraph*{#1}\mbox{}}
  \newcommand{\xxxParagraphNoStar}[1]{\oldparagraph{#1}\mbox{}}
\fi
\ifx\subparagraph\undefined\else
  \let\oldsubparagraph\subparagraph
  \renewcommand{\subparagraph}{
    \@ifstar
      \xxxSubParagraphStar
      \xxxSubParagraphNoStar
  }
  \newcommand{\xxxSubParagraphStar}[1]{\oldsubparagraph*{#1}\mbox{}}
  \newcommand{\xxxSubParagraphNoStar}[1]{\oldsubparagraph{#1}\mbox{}}
\fi
\makeatother


\usepackage{longtable,booktabs,array}
\newcounter{none} % for unnumbered tables
\usepackage{calc} % for calculating minipage widths
% Correct order of tables after \paragraph or \subparagraph
\usepackage{etoolbox}
\makeatletter
\patchcmd\longtable{\par}{\if@noskipsec\mbox{}\fi\par}{}{}
\makeatother
% Allow footnotes in longtable head/foot
\IfFileExists{footnotehyper.sty}{\usepackage{footnotehyper}}{\usepackage{footnote}}
\makesavenoteenv{longtable}
\usepackage{graphicx}
\makeatletter
\newsavebox\pandoc@box
\newcommand*\pandocbounded[1]{% scales image to fit in text height/width
  \sbox\pandoc@box{#1}%
  \Gscale@div\@tempa{\textheight}{\dimexpr\ht\pandoc@box+\dp\pandoc@box\relax}%
  \Gscale@div\@tempb{\linewidth}{\wd\pandoc@box}%
  \ifdim\@tempb\p@<\@tempa\p@\let\@tempa\@tempb\fi% select the smaller of both
  \ifdim\@tempa\p@<\p@\scalebox{\@tempa}{\usebox\pandoc@box}%
  \else\usebox{\pandoc@box}%
  \fi%
}
% Set default figure placement to htbp
\def\fps@figure{htbp}
\makeatother



\ifLuaTeX
\usepackage[bidi=basic,shorthands=off]{babel}
\else
\usepackage[bidi=default,shorthands=off]{babel}
\fi
\ifLuaTeX
  \usepackage{selnolig} % disable illegal ligatures
\fi


\setlength{\emergencystretch}{3em} % prevent overfull lines

\providecommand{\tightlist}{%
  \setlength{\itemsep}{0pt}\setlength{\parskip}{0pt}}




\usepackage[]{natbib}
\bibliographystyle{elsarticle-harv}


\usepackage{orcidlink}
\usepackage{float}
\makeatletter
\@ifpackageloaded{caption}{}{\usepackage{caption}}
\AtBeginDocument{%
\ifdefined\contentsname
  \renewcommand*\contentsname{Table of contents}
\else
  \newcommand\contentsname{Table of contents}
\fi
\ifdefined\listfigurename
  \renewcommand*\listfigurename{List of Figures}
\else
  \newcommand\listfigurename{List of Figures}
\fi
\ifdefined\listtablename
  \renewcommand*\listtablename{List of Tables}
\else
  \newcommand\listtablename{List of Tables}
\fi
\ifdefined\figurename
  \renewcommand*\figurename{Figure}
\else
  \newcommand\figurename{Figure}
\fi
\ifdefined\tablename
  \renewcommand*\tablename{Table}
\else
  \newcommand\tablename{Table}
\fi
}
\@ifpackageloaded{float}{}{\usepackage{float}}
\floatstyle{ruled}
\@ifundefined{c@chapter}{\newfloat{codelisting}{h}{lop}}{\newfloat{codelisting}{h}{lop}[chapter]}
\floatname{codelisting}{Listing}
\newcommand*\listoflistings{\listof{codelisting}{List of Listings}}
\makeatother
\makeatletter
\makeatother
\makeatletter
\@ifpackageloaded{caption}{}{\usepackage{caption}}
\@ifpackageloaded{subcaption}{}{\usepackage{subcaption}}
\makeatother
\journal{Computers and Electronics in Agriculture}
\usepackage{bookmark}
\IfFileExists{xurl.sty}{\usepackage{xurl}}{} % add URL line breaks if available
\urlstyle{same}
\hypersetup{
  pdftitle={Computational Scalability of Deterministic and Two-Stage Stochastic MILP Models for Grain-Storage Network Design and Hub Location},
  pdfauthor={Victor Rafael Rezende Celestino; Artur Guerra Rosa; Andréia Elizabeth Silva Barros; Gabriela Corsano; Luis Javier Zeballos; Rodolfo Dondo; Silvia Araujo dos Reis},
  pdflang={en},
  pdfkeywords={grain storage, network design, hub
location, mixed-integer linear programming, two-stage stochastic
programming, computational scalability, high-performance computing},
  colorlinks=true,
  linkcolor={blue},
  filecolor={Maroon},
  citecolor={Blue},
  urlcolor={Blue},
  pdfcreator={LaTeX via pandoc}}


\setlength{\parindent}{6pt}
\begin{document}

\begin{frontmatter}
\title{Computational Scalability of Deterministic and Two-Stage
Stochastic MILP Models for Grain-Storage Network Design and Hub
Location}
\author[1]{Victor Rafael Rezende
Celestino\,\orcidlink{0000-0001-5913-2997}%
\corref{cor1}%
}
 \ead{vrcelestino@unb.br}
\author[1]{Artur Guerra Rosa\,\orcidlink{0000-0002-5013-4408}%
%
}
 \ead{artur.guerra@unb.br}
\author[1]{Andréia Elizabeth Silva
Barros\,\orcidlink{0000-0002-7409-0550}%
%
}
 \ead{andreiabarros@tc.df.gov.br}
\author[2]{Gabriela Corsano\,\orcidlink{0000-0002-0235-4262}%
%
}
 \ead{gcorsano@santafe-conicet.gov.ar}
\author[2]{Luis Javier Zeballos\,\orcidlink{0000-0002-8548-2912}%
%
}
 \ead{zeballos@intec.unl.edu.ar}
\author[2]{Rodolfo Dondo\,\orcidlink{0000-0003-4580-0493}%
%
}
 \ead{rdondo@santafe-conicet.gov.ar}
\author[1]{Silvia Araujo dos Reis\,\orcidlink{0000-0002-1646-4454}%
%
}
 \ead{silviareis@unb.br}

\affiliation[1]{organization={University of
Brasilia},city={Brasilia},country={Brazil},countrysep={,},postcodesep={}}
\affiliation[2]{organization={National University of the
Litoral},city={Santa
Fe},country={Argentina},countrysep={,},postcodesep={}}

\cortext[cor1]{Corresponding author}








\begin{abstract}
Brazilian grain production reached 352 million metric tons in the 2024/25 crop year, while national static storage capacity remains below 65\% of that volume and logistics depend on a non-optimized network. Brazilian government has investment plans for high performance computing (HPC) infrastructure with the intent to prioritize agricultural research.
Designing grain-storage networks requires decisions about facility location, storage capacity and grain handling before uncertain supply and demand are known. The resulting mixed-integer linear programming (MILP) can contain millions of strategic decisions, making computational requirements central to their use in agricultural planning. This study evaluates deterministic and two-stage stochastic models through computing upscale, speedup, efficiency and scaling assessment of MILP solution in HPC. Warehouse-only instances contain 215, 300, 400 and 500 facilities, with nine supply--demand scenarios in the stochastic model. The stochastic formulation grows from 14.1 to 63.4 million variables across this range. All nine warehouse-only observations satisfy the prescribed 1\% relative-gap criterion; the larger 500-facility instance reaches a 0.21\% gap in approximately 128 minutes of native optimization. In the common Pyomo h300 experiment, a commercial solver (Gurobi) reached the 1\% target in all 15 runs and a open-source solver (HiGHS) reached it in three one-thread runs. The returned stochastic plans allocate capacity investment across the admitted candidate sites, indicating distributed capacity sizing within the available network. However, substantial handling and storage exceedances and unmet demand are necessary to maintain solutions feasibility. The first partial h300 root-method experiment additionally suggests that the initial LP relaxation can dominate solver effort even when the branch-and-bound tree contains few explored nodes. These preliminary results motivate a controlled comparison of root algorithms, incumbent heuristics and physical-core allocation, without establishing a general speedup. Evaluating solution quality together with computation metrics provide decision support about optimized plans for grain-storage planning.
\end{abstract}


\begin{keyword}
    grain storage \sep network design \sep hub
location \sep mixed-integer linear programming \sep two-stage stochastic
programming \sep computational scalability \sep
    high-performance computing
\end{keyword}
\end{frontmatter}


\section{Introduction}\label{introduction}
Brazilian agribusiness generates nearly a quarter of the national gross domestic product and accounts for over 40 percent of national exports \citep{basso_agriculture_2024}. Supported by technological intensification and total factor productivity growth over several decades, national grain production reached 352 million metric tons in the 2024/25 crop year, mainly due to soybean and maize yields destined for both domestic and international markets \citep{conab_portal_2026b,mascarenhas2014}. However, national static storage capacity remains below 65\% of that volume \citep{conab_portal_2026a} and logistics depend on a non-optimized network.

Brazil's Artificial Intelligence Plan identifies
high-performance computing infrastructure and technical capacity as priorities for research and innovation \citep{mcti2025pbia}. The National Laboratory for Scientific Computing is developing a separate AI-oriented supercomputing installation in Macaíba, Rio Grande do Norte \citep{lncc2026supercomputador}. Research in agricultural operational and logistic efficiencies are included as priority topics to be addressed from infrastructure investments \citep{pera2026}.
Embrapa's AgroDigital initiative has acquired high-performance computing and GPU servers to support research that requires substantial data
processing and storage \citep{embrapa2024datacenter}. These initiatives illustrate the importance of high performance computing (HPC) infrastructure across this research domain.

Grain-storage networks connect seasonal harvests with demand distributed
across places and periods. Choosing where to open or expand warehouses
affects how much grain can be held, how quickly it can be received and
dispatched, and which markets the network can serve.
Grain-storage planning raises a specific computational problem. The
selection of warehouse sites and the sizing of storage and handling
capacity combine discrete investment choices with flows that vary across
products, periods and scenarios. A larger network increases the number
of decisions and can place substantial demands on model construction,
memory and solution time. This study uses experiments at an HPC infrastruture to
examine the conditions under which deterministic and two-stage
stochastic mixed-integer linear programs yield optimization plans for soybean and corn storage networks. The analysis relates measured computation metrics and attainable solution quality to the investment and strategic decisions that the models
return.

Grain-storage location and integrated network design are established research topics \citep{essien2018, gholamian2017}. Agricultural network
uncertainty has also been addressed through wheat-network design and tactical soybean planning \citep{hosseini2020, reis2023}. Decomposition
and distributed computing already provide computational precedents in agricultural MILP \citep{shastri2011}. The study connects the effort
needed to construct and solve these models with the certified quality and agricultural implications of the plans they return.

Therefore, the research question is: How computational scalability affects
the optimal resolution of grain-storage network design and hub location models? Optimality is evaluated in terms of MILP solution quality attainable.

A complementary diagnostic question concerns whether the initial LP relaxation, rather than extensive branch-and-bound exploration, accounts for a substantial share of solution effort, and how its observed time and incumbent attainment vary with the root method, heuristic profile and physical-core allocation.

The general objective is to investigate the effects of computational scalability in the resolution of grain-storage network design and hub
location models in an HPC infrastruture.
The specific objectives seek three distinct contributions. C1 characterizes observed
matrix dimensions, construction and solution effort, memory and certified gaps across the evaluated populations. C2 interprets the returned investment and operating decisions by separating financial
costs, artificial penalties, domestic service and capacity exceedances.
C3 documents human-supervised AI assistance and verification as a complementary account of research practice, detailed in Appendix B.
Table~\ref{tbl-objectives} links the four specific
objectives to the contributions they support.

\begin{longtable}[]{@{}
  >{\raggedright\arraybackslash}p{(\linewidth - 4\tabcolsep) * \real{0.3333}}
  >{\raggedright\arraybackslash}p{(\linewidth - 4\tabcolsep) * \real{0.3333}}
  >{\raggedright\arraybackslash}p{(\linewidth - 4\tabcolsep) * \real{0.3333}}@{}}
\caption{Research objectives.}\label{tbl-objectives}\tabularnewline
\toprule\noalign{}
\begin{minipage}[b]{\linewidth}\raggedright
Specific Objective
\end{minipage} & \begin{minipage}[b]{\linewidth}\raggedright
Scope
\end{minipage} & \begin{minipage}[b]{\linewidth}\raggedright
Contribution
\end{minipage} \\
\midrule\noalign{}
\endfirsthead
\toprule\noalign{}
\begin{minipage}[b]{\linewidth}\raggedright
Objective
\end{minipage} & \begin{minipage}[b]{\linewidth}\raggedright
Scope
\end{minipage} & \begin{minipage}[b]{\linewidth}\raggedright
Contribution
\end{minipage} \\
\midrule\noalign{}
\endhead
\bottomrule\noalign{}
\endlastfoot
SO1 Mathematical modeling & Implement deterministic and two-stage
stochastic MILPs for warehouse location, capacity investment and
grain-network operation & C1, C2 \\
SO2 Computational scalability & Measure model construction, memory,
optimization time, allocated resources and attainable solution quality
as warehouse populations increase; additionally characterize root-relaxation effort in an exploratory h300 experiment & C1 \\
SO3 Agricultural interpretation & Interpret validated location,
capacity-investment and grain-distribution decisions across h215, h300,
h400 and h500  & C2 \\
SO4 Complementary methodology & Document human-supervised AI assistance,
scientific decision provenance, verification and reproducibility & C3 \\
\end{longtable}

The next section situates these objectives in agricultural planning and computational optimization. The methods and results then connect the
model, resource measurements and independently checked solutions to storage-network decisions.


\section{Related work and research
positioning}\label{related-work-and-research-positioning}

\subsection{Grain-storage planning as an integrated decision
problem}\label{grain-storage-planning-as-an-integrated-decision-problem}

The systematic review by \citet{rosa2025} places grain-warehouse
location within a broader family of mathematical planning models.
Complementarily, \citet{garima2024} organize grain supply-chain decision
research around operational complexity, risk, infrastructure,
stakeholders and solution methods. Together, these perspectives show why
adding sites to a location model is insufficient: a strategic design
must connect storage, seasonal handling and access to markets. The
following synthesis compares published approaches along these
substantive dimensions.
\href{https://unb-lamfo-or-ai-hpc.github.io/model-agrologistic/supplementary-review.html}{Appendix
A} provides the search strategy, selection record and the complete systematic literature review.

Integrated wheat-network design in \citet{gholamian2017} links
infrastructure choices to the organization of the supply chain.
Scenario-based wheat planning in \citet{pourmohammadi2020} and
sustainable resilient network design in \citet{hosseini2020} extend
this perspective to uncertain conditions and competing performance
concerns. These contributions motivate an integrated treatment, but
their commodity systems, constraints and objectives cannot be
transferred unchanged to a Brazilian grain-storage network. In
particular, resilience to disruption, seasonal variation and a numerical
relaxation of capacity are different concepts. A slack preserves a
mathematical solution; it does not by itself provide the transport,
temporary equipment or institutional arrangements that would make the
physical network resilient.

\citet{essien2018} address stakeholder objectives in grain-storage
planning, and \citet{filippi2022} examine the governance of rural
warehouse condominiums. In the goal-programming model of
\citet{essien2018}, the number of sites is prescribed and stakeholder
transport targets enter the objective. Here the model sizes capacity and
links investment to scenario-dependent operation. Land access, financing
and operating agreements remain part of the appraisal of any candidate
location.

\subsection{Closest precedents and the scope of the empirical
question}\label{closest-precedents-and-the-scope-of-the-empirical-question}

The closest agricultural studies establish why location and uncertainty alone cannot define this article's gap.
Table~\ref{tbl-closest-precedents} distinguishes their decision or computational emphasis from the assessment undertaken here. The studies differ in decision horizon, infrastructure choices and solution method.
These distinctions determine which aspects can inform the present
assessment.

\begin{longtable}[]{@{}
  >{\raggedright\arraybackslash}p{(\linewidth - 4\tabcolsep) * \real{0.3333}}
  >{\raggedright\arraybackslash}p{(\linewidth - 4\tabcolsep) * \real{0.3333}}
  >{\raggedright\arraybackslash}p{(\linewidth - 4\tabcolsep) * \real{0.3333}}@{}}
\caption{Closest agricultural and computational precedents}\label{tbl-closest-precedents}\tabularnewline
\toprule\noalign{}
\begin{minipage}[b]{\linewidth}\raggedright
Precedent
\end{minipage} & \begin{minipage}[b]{\linewidth}\raggedright
Established decision or computational emphasis
\end{minipage} & \begin{minipage}[b]{\linewidth}\raggedright
Research question examined
\end{minipage} \\
\midrule\noalign{}
\endfirsthead
\toprule\noalign{}
\begin{minipage}[b]{\linewidth}\raggedright
Precedent
\end{minipage} & \begin{minipage}[b]{\linewidth}\raggedright
Established decision or computational emphasis
\end{minipage} & \begin{minipage}[b]{\linewidth}\raggedright
Question examined here
\end{minipage} \\
\midrule\noalign{}
\endhead
\bottomrule\noalign{}
\endlastfoot
\citet{essien2018} & Multi-stakeholder grain-storage location and
decision support & How do returned investment and service outcomes
relate to a certified scalar objective? \\
\citet{gholamian2017} & Integrated wheat-network planning, including
silo location and distribution & What construction, solution and memory
requirements accompany explicitly defined extensive-form instances? \\
\citet{pourmohammadi2020} & Scenario-based location and discrete storage
sizing with hard demand fulfillment & How do continuous capacity sizing
and penalized shortages change the interpretation of returned plans? \\
\citet{hosseini2020} & Wheat-network design/redesign with sustainability
and resilience concerns & What do the present cost and slack variables
support, without adding unmodeled resilience or sustainability
outcomes? \\
\citet{reis2023} & Two-stage tactical soybean LP with purchasing and
transport contracts; 243 scenarios & How does shared infrastructure
investment differ from tactical contracting under uncertainty? \\
\citet{shastri2011} & Decomposition and distributed solution of an
agricultural MILP & What attainable quality and resource use are
observed for the present direct extensive-form solution? \\
\end{longtable}

These studies motivate distinct comparisons. Discrete sizing and hard
demand fulfillment in \citet{pourmohammadi2020} clarify how continuous
capacity choices and penalized shortage affect the interpretation of a
returned plan. \citet{hosseini2020} treat sustainability and resilience
as modeled concerns. \citet{shastri2011} evaluate
decomposition and distributed solution.
This study asks a different planning question and adds explicit attention to construction, attained integer-solution quality, computation resorces use and the physical interpretation of the returned capacity plan for grain-storage strategic decisions.

\subsection{Network representation and decision
horizons}\label{network-representation-and-decision-horizons}

Dynamic-network models connect storage investment to grain flows over
time \citep{foulds2005}. In a Brazilian soybean setting,
\citet{mascarenhas2014} study warehousing location with road and
intermodal export connections. The distinction is consequential: a
throughput facility that transfers arrivals onward is not
interchangeable with a warehouse that carries inventory between months.
Static stock capacity, annual throughput and vehicle capacity must
therefore remain separate dimensions in a benchmark.
The logistics-integration-center formulation of \citet{guimaraes2017}
preserves origin-destination demand relationships. By contrast, the
present model conserves each product but does not retain an origin label
after it enters a warehouse. Product from different origins can satisfy
the same customer's demand. This assumption is appropriate for a
fungible-commodity planning abstraction, but not for tracing contracts,
quality lots or individual consignments.

Terminal siting in \citet{morais2023} connects location, capacity
choices and access to railways or waterways. \citet{milanez2026}
integrate soybean export network design with hub functions and economies
of scale. One variant permits seaport-capacity overflow at a penalty, reporting CPU time and optimality gaps. Penalized capacity relief
and computational quality reporting therefore have direct precedents.
These studies motivate reporting storage and transport decisions
together. Multimodal models also represent modal
transfers and transport beyond those outlets.
The present network represents road movements to domestic destinations and export ports.
Table~\ref{tbl-literature} summarizes the selected agricultural network-design studies, their
modeling distinctions and implications for comparbility.

\begin{longtable}[]{@{}
  >{\raggedright\arraybackslash}p{(\linewidth - 4\tabcolsep) * \real{0.3333}}
  >{\raggedright\arraybackslash}p{(\linewidth - 4\tabcolsep) * \real{0.3333}}
  >{\raggedright\arraybackslash}p{(\linewidth - 4\tabcolsep) * \real{0.3333}}@{}}
\caption{Selected agricultural network-design studies}\label{tbl-literature}\tabularnewline
\toprule\noalign{}
\begin{minipage}[b]{\linewidth}\raggedright
Study
\end{minipage} & \begin{minipage}[b]{\linewidth}\raggedright
Network and decision emphasis
\end{minipage} & \begin{minipage}[b]{\linewidth}\raggedright
Comparison implication
\end{minipage} \\
\midrule\noalign{}
\endfirsthead
\toprule\noalign{}
\begin{minipage}[b]{\linewidth}\raggedright
Study
\end{minipage} & \begin{minipage}[b]{\linewidth}\raggedright
Network and decision emphasis
\end{minipage} & \begin{minipage}[b]{\linewidth}\raggedright
Comparison implication
\end{minipage} \\
\midrule\noalign{}
\endhead
\bottomrule\noalign{}
\endlastfoot
\citet{foulds2005} & Dynamic grain-silo location and flow & Separate
network size, capacity decisions and temporal flow \\
\citet{mascarenhas2014} & Mato Grosso soybean warehouses and export
connections & Distinguish throughput facilities from stock-carrying
warehouses \\
\citet{guimaraes2017} & Logistics integration centers and peer-to-peer
demand & Commodity mixing differs from origin-destination
preservation \\
\citet{morais2023} & Capacitated terminals with intermodal location
preferences & Modal access is not equivalent to a road-distance proxy \\
\citet{milanez2026} & Hub location and soybean-export network design &
Report hub function, capacity and route use together \\
\citet{granillo2019} & Contract-farming network with simulation and MILP
& Simulation-based uncertainty is not automatically two-stage
recourse \\
\citet{mendoza2021} & Scenario-based agro-food facility location & State
which decisions precede the uncertain outcome \\
\citet{vazquez2021} & Guayule/guar supply-chain design with
sustainability criteria & Processing yields and sustainability
constraints change comparability \\
\citet{mogale2020} & Multi-echelon grain network with risk and losses &
Losses and discrete transport resources require additional balances \\
\end{longtable}

\subsection{Parameters, decisions and objective
functions}\label{parameters-decisions-and-objective-functions}

Economic formulations typically connect fixed facility choices to
continuous material flows. The opening variable controls whether
infrastructure can be used, while installed capacity determines how much
can be stored or handled. A benchmark must distinguish a fixed-size site
choice from scalable capacity: both may contain a binary opening
variable, but their investment feasible sets differ. Counting producing
regions or candidate sites is not an integer decision variable unless
the optimizer actually selects that count.
Agricultural models also differ in the economic boundary. While some consider freight and
construction alone, others include freight, inventory, processing, emissions, losses and employment in the objective.
\citet{mogale2020} include risk and post-harvest losses in a green food
supply-chain design. \citet{vazquez2021} connect supply-chain design to
environmental and social impacts; \citet{mogale2023} examine social
sustainability. Those dimensions are relevant to public planning but are
outside the scope of the present formulation.

Logistics costs depend on more than distance.
\citet{kussano2012} analyze soybean export logistics using a broader
cost accounting boundary, including factors that can alter the ranking
of modal alternatives. Accordingly, road-network distances associate freight approximation based on modeled loads and unit transport costs.
Congestion, taxes, service contracts and seasonally changing tariffs
remain possible sources of divergence between the least-cost modeled
route and actual grain movements.
Spatial suitability and network optimization are complementary rather
than substitutable. \citet{souza2024} use spatial multicriteria analysis
for agricultural warehouse location; such screening can constrain the
candidate set before a flow model allocates capacity. The present
candidate population preserves geographic and facility-type diversity
but does not replace environmental, land-use or engineering suitability
assessment. Similarly, the wider agricultural and carbon implications of
Brazilian transport infrastructure examined by \citet{wang2024} lie
outside a freight-and-storage objective. A shorter modeled route cannot
be interpreted automatically as a net environmental gain.

\subsection{Uncertainty and strategic
interpretation}\label{uncertainty-and-strategic-interpretation}

Uncertainty can enter a model through simulation, scenarios, robust
constraints or probabilistic objectives. These approaches answer
different questions. \citet{granillo2019} combine a deterministic
optimization model with simulation in a barley contract-farming network.
\citet{mendoza2021} study scenario-based facility location, while the
adoption scenarios in \citet{vazquez2021} concern a different source of
uncertainty.
\citet{pourmohammadi2020} provide a closer stochastic infrastructure
precedent: six equiprobable supply-demand scenarios, shared choices
among four facility-capacity levels, and scenario-dependent operation.
Their assumptions prohibit shortage and include imports, quality
blending and road/rail transport. The present research instead sizes
capacity continuously, represents road movements and permits penalized
shortages. These differences change the feasible set and the meaning of
optimal cost outcomes.
\citet{maiyar2019} model emergency hubs as physical options for
disruptions. The slacks used in the present study quantify excess demand
on nominal capacity within the modeled network.
Integrated decision support links planning horizons and optimization \citep{mardaneh2021}. The joint
crop-establishment, harvesting and network-design perspective of \citet{anselmino2025} provides another example, in bioenergy rather than
grain storage. Such work motivates clear interfaces between data preparation, strategic investment and network evaluation.

\subsection{Computational approaches in agricultural
systems}\label{computational-approaches-in-agricultural-systems}

\citet{shastri2011} address a large agricultural optimization model
through a decomposition and distributed-computing approach to BioFeed, a
biomass supply planning problem. This is direct precedent for treating
computational solution strategy, rather than the use of a commercial
optimizer alone, as a research contribution.
\citet{knapen2025} use standard big-data and distributed-computing
technologies for scalable crop-growth simulations. Their application
illustrates a different parallelization opportunity: distributing
simulation work is not equivalent to solving a MILP whose scenarios
share first-stage integer decisions. \citet{zhao2012} study
parallelization and optimization of spatial analysis for environmental
model-data assembly. This distinction supports measuring data
preparation separately from the numerical solver rather than attributing
all elapsed time to optimization.
Table~\ref{tbl-computational-literature} relates the agricultural
planning problem to computational approaches at the optimization,
simulation and data-processing levels.

\begin{longtable}[]{@{}
  >{\raggedright\arraybackslash}p{(\linewidth - 4\tabcolsep) * \real{0.3333}}
  >{\raggedright\arraybackslash}p{(\linewidth - 4\tabcolsep) * \real{0.3333}}
  >{\raggedright\arraybackslash}p{(\linewidth - 4\tabcolsep) * \real{0.3333}}@{}}
\caption{Computational approaches for agricultural planning}\label{tbl-computational-literature}\tabularnewline
\toprule\noalign{}
\begin{minipage}[b]{\linewidth}\raggedright
Computational layer
\end{minipage} & \begin{minipage}[b]{\linewidth}\raggedright
Published precedent
\end{minipage} & \begin{minipage}[b]{\linewidth}\raggedright
Present implementation and evidence boundary
\end{minipage} \\
\midrule\noalign{}
\endfirsthead
\toprule\noalign{}
\begin{minipage}[b]{\linewidth}\raggedright
Computational layer
\end{minipage} & \begin{minipage}[b]{\linewidth}\raggedright
Published precedent
\end{minipage} & \begin{minipage}[b]{\linewidth}\raggedright
Present implementation and evidence boundary
\end{minipage} \\
\midrule\noalign{}
\endhead
\bottomrule\noalign{}
\endlastfoot
Agricultural MILP solution strategy & \citet{shastri2011} & Direct
extensive-form solution; decomposition is not implemented in the
evaluated scalar pipeline \\
Distributed agricultural simulation & \citet{knapen2025} & Slurm
experiment scheduling, not a distributed scenario solver \\
Spatial data preparation & \citet{zhao2012} & Versioned OSRM
materialization and pair-cache reuse; no matched speedup benchmark \\

\end{longtable}

\noindent The large stochastic extensive form also invites a distinction between the cost of solving the initial linear-programming (LP) relaxation and the work performed after branching. A single explored branch-and-bound node does not imply a computationally inexpensive solve, nor prove that the model structurally prevents branching. This distinction is examined below using recorded solver logs and an explicitly exploratory experiment.

\section{Materials and methods}\label{materials-and-methods}

The investment formulation follows \citet{rosa2026thesis} in its
treatment of facility opening, expansion and bulk-handling conversion.
Bulk-handling conversion is usually implemented with additional equipment and process improvements rather than capacity expansion.
The objective function accounts for investments once and weights recourse costs by scenario probability. Investment decisions are shared across
scenarios, and each scenario has a separate operating trajectory
\citep{birge2011}.

\subsection{Planning problem and
notation}\label{planning-problem-and-notation}

The problem selects a grain-storage network and allocates exogenous
supply over a finite planning horizon. Existing facilities are
available. Candidate opening, candidate capacity, expansion and
bulk-handling investment precede the first period. Logistics and
permitted shortages/capacity exceedances depend on the realized
supply/demand trajectory. The model conserves each product; it omits
production choices, quality loss, processing, transport lead times and
fleet scheduling. Domestic demand has a shortage variable; exports are
optional outlets with their declared upper bounds. All supply must leave
its origin. Initial inventory is supplied, and terminal inventory is
free within the constraints.

Table~\ref{tbl-sets} defines the sets, Table~\ref{tbl-parameters} gives
the coefficients and their units, and Table~\ref{tbl-variables} gives
the decisions and their domains. Flows are defined on admissible arcs.
Ineligible investment quantities are zero. Each scenario satisfies the
physical balances and capacity limits. Scenario probabilities weight the
recourse costs.

The arc families connect origins to warehouses, warehouses to
destinations, and warehouses to other warehouses. Configuration D also
allows direct origin-to-destination movements. The warehouse-only
sequence uses configuration W. Appendix C gives the exact crosswalk
between these families and the symbols in the equations.

\begin{longtable}[]{@{}
  >{\raggedright\arraybackslash}p{(\linewidth - 2\tabcolsep) * \real{0.5000}}
  >{\raggedright\arraybackslash}p{(\linewidth - 2\tabcolsep) * \real{0.5000}}@{}}
\caption{Model sets and index domains.}\label{tbl-sets}\tabularnewline
\toprule\noalign{}
\begin{minipage}[b]{\linewidth}\raggedright
Symbol
\end{minipage} & \begin{minipage}[b]{\linewidth}\raggedright
Definition
\end{minipage} \\
\midrule\noalign{}
\endfirsthead
\toprule\noalign{}
\begin{minipage}[b]{\linewidth}\raggedright
Symbol
\end{minipage} & \begin{minipage}[b]{\linewidth}\raggedright
Definition
\end{minipage} \\
\midrule\noalign{}
\endhead
\bottomrule\noalign{}
\endlastfoot
\(O\); \(W=W^E\mathbin{\dot\cup}W^N\) & Supply origins; existing and
candidate warehouses \\
\(W^X,W^B\subseteq W\) & Sites eligible for expansion and bulk-handling
conversion \\
\(C^D,C^X\) & Domestic and export destinations, retaining distinct roles
at the same municipality \\
\(P\); \(T\); \(S\) & Products; ordered finite planning periods;
scenarios \\
\(A^{OD},A^{DC},A^{DD},A^{OC}\) & Product-specific directed
origin--warehouse, warehouse--customer, interwarehouse and optional
direct origin--customer arcs \\
\end{longtable}

\begin{longtable}[]{@{}
  >{\raggedright\arraybackslash}p{(\linewidth - 2\tabcolsep) * \real{0.5000}}
  >{\raggedright\arraybackslash}p{(\linewidth - 2\tabcolsep) * \real{0.5000}}@{}}
\caption{Model coefficients and
dimensions.}\label{tbl-parameters}\tabularnewline
\toprule\noalign{}
\begin{minipage}[b]{\linewidth}\raggedright
Symbol
\end{minipage} & \begin{minipage}[b]{\linewidth}\raggedright
Definition and units
\end{minipage} \\
\midrule\noalign{}
\endfirsthead
\toprule\noalign{}
\begin{minipage}[b]{\linewidth}\raggedright
Symbol
\end{minipage} & \begin{minipage}[b]{\linewidth}\raggedright
Definition and units
\end{minipage} \\
\midrule\noalign{}
\endhead
\bottomrule\noalign{}
\endlastfoot
\(a_{opt}\); \(D_{cpt}\); \(E_{cpt}\) & Supply, domestic demand and
export upper bound, t within a period \\
\(I^0_{wp}\); \(K^0_w\) & Initial stock and initial static capacity,
t \\
\(R^0_w,H^0_w\); \(d_t\) & Initial reception/shipping rates, t/day;
operating days per period \\
\(\overline k_w,\overline e_w,\overline b_w\) & Candidate, expansion and
conversion upper bounds, t on the respective investment basis \\
\(\rho^N,\rho^X,\rho^B\); \(\eta^N,\eta^X,\eta^B\) & Daily
reception/shipping gains per invested tonne, day\(^{-1}\) \\
\(\ell^{OD}_{ow},\ell^{DC}_{wc},\ell^{DD}_{wv},\ell^{OC}_{oc}\) &
Directed road distances, km \\
\(f^O_o,f^W_w\); \(\alpha\); \(\tau_w\) & Freight rates, BRL/(t km);
dimensionless interwarehouse multiplier; receiving charge, BRL/t \\
\(h_{wp}\) & Effective charge for holding one tonne for one modeled
period, BRL/t; the source tariff is quoted in BRL/(t period) \\
\(F^N_w,F^X_w,F^B_w\) & Fixed opening, expansion and conversion charges,
BRL \\
\(c^N_w,c^X_w,c^B_w\) & Corresponding quantity-dependent investment
charges, BRL per investment-tonne \\
\(\lambda^D_{cp},\lambda^K_w,\lambda^R_w\) & Shortage, stock-exceedance
and reception-overflow weights; BRL per corresponding slack tonne at its
period \\
\(\pi_s>0\), \(\sum_{s\in S}\pi_s=1\) & Scenario probabilities; the
present STO experiments use nine weights of \(1/9\) \\
\end{longtable}

\begin{longtable}[]{@{}
  >{\raggedright\arraybackslash}p{(\linewidth - 2\tabcolsep) * \real{0.5000}}
  >{\raggedright\arraybackslash}p{(\linewidth - 2\tabcolsep) * \real{0.5000}}@{}}
\caption{Decisions and derived
expressions.}\label{tbl-variables}\tabularnewline
\toprule\noalign{}
\begin{minipage}[b]{\linewidth}\raggedright
Symbol
\end{minipage} & \begin{minipage}[b]{\linewidth}\raggedright
Meaning, domain and units
\end{minipage} \\
\midrule\noalign{}
\endfirsthead
\toprule\noalign{}
\begin{minipage}[b]{\linewidth}\raggedright
Symbol
\end{minipage} & \begin{minipage}[b]{\linewidth}\raggedright
Meaning, domain and units
\end{minipage} \\
\midrule\noalign{}
\endhead
\bottomrule\noalign{}
\endlastfoot
\(X=(y,z,g,k,e,b)\) & Shared first-stage investment vector \\
\(y_w,z_w,g_w\) & Binary candidate opening, expansion and conversion, on
\(W^N,W^X,W^B\) \\
\(k_w,e_w,b_w\ge0\) & Corresponding investment quantities, t on their
declared basis \\
\(x^{OD}_{owpt},x^{DC}_{wcpt},x^{DD}_{wvpt},x^{OC}_{ocpt}\ge0\) & Grain
flow on instantiated arcs, t within the period \\
\(I_{wpt},u_{cpt},q^K_{wt},q^R_{wt}\ge0\) & End stock, domestic
shortage, stock exceedance and reception overflow, t at/within the
period \\
\(A_w\) & Activation expression: \(1\) for \(w\in W^E\), \(y_w\) for
\(w\in W^N\) \\
\(K_w(X),R_w(X),H_w(X)\) & Derived static t and daily reception/shipping
t/day capacities \\
\(B^{in}_{wpt},B^{out}_{wpt}\) & Derived incoming/outgoing flow sums, t
within the period \\
\(Y\); \(Y^s\) & DET operating vector; scenario-\(s\) flow, stock and
slack vector \\
\end{longtable}

Candidates have zero initial static and handling capacity. Opening a
candidate or expanding an existing warehouse can add static capacity,
measured in tonnes of grain that can be held at a given time.
Bulk-handling conversion raises the daily reception and shipping rates
of an eligible warehouse. These rates are expressed in tonnes per day
and converted to period limits. A warehouse can therefore handle more
arrivals and departures while retaining the same limit on grain held in
storage. The model makes these investments once, before the first
planning period.

\subsection{Deterministic model (DET)}\label{deterministic-model-det}

Costs are undiscounted BRL totals over the horizon. Receiving handling
is part of the OD/DD unit transport/handling cost and is charged once at
receipt. Per-period holding and stock-exceedance coefficients charge one
modeled period at each time index; that unit-period duration is implicit
in D1. Daily handling rates instead require the explicit operating-day
factor in the capacity rows. Appendix D provides the dimensional
crosswalk and worked conversions.

\[
\begin{aligned}
\min_{X,Y}\ Z_D={}&C^{INV}(X)+C^{TR}(Y)+C^{HOLD}(Y)+C^{PEN}(Y), &&\text{(D1)}\\
C^{INV}(X)={}&\sum_{w\in W^N}(F^N_wy_w+c^N_wk_w)
 +\sum_{w\in W^X}(F^X_wz_w+c^X_we_w)
 +\sum_{w\in W^B}(F^B_wg_w+c^B_wb_w),\\
C^{TR}(Y)={}&\sum_{t\in T}\sum_{(o,w,p)\in A^{OD}}
 (f^O_o\ell^{OD}_{ow}+\tau_w)x^{OD}_{owpt}\\
&+\sum_{t\in T}\sum_{(w,c,p)\in A^{DC}}f^W_w\ell^{DC}_{wc}x^{DC}_{wcpt}\\
&+\sum_{t\in T}\sum_{(w,v,p)\in A^{DD}}
 (\alpha f^W_w\ell^{DD}_{wv}+\tau_v)x^{DD}_{wvpt}\\
&+\sum_{t\in T}\sum_{(o,c,p)\in A^{OC}}f^O_o\ell^{OC}_{oc}x^{OC}_{ocpt},\\
C^{HOLD}(Y)={}&\sum_{w\in W,p\in P,t\in T}h_{wp}I_{wpt},\\
C^{PEN}(Y)={}&\sum_{c\in C^D,p\in P,t\in T}\lambda^D_{cp}u_{cpt}
 +\sum_{w\in W,t\in T}(\lambda^K_wq^K_{wt}+\lambda^R_wq^R_{wt}).
\end{aligned}
\]

\[
\begin{aligned}
0\le k_w&\le\overline k_wy_w &&\forall w\in W^N, &&\text{(D2)}\\
0\le e_w&\le\overline e_wz_w,\quad z_w\le A_w &&\forall w\in W^X,\\
0\le b_w&\le\overline b_wg_w,\quad g_w\le A_w &&\forall w\in W^B,\\
z_w+g_w&\le A_w &&\forall w\in W^X\cap W^B,\\
K_w(X)&=K^0_w+k_w+e_w &&\forall w\in W, &&\text{(D3)}\\
R_w(X)&=R^0_w+\rho^Nk_w+\rho^Xe_w+\rho^Bb_w &&\forall w\in W,\\
H_w(X)&=H^0_w+\eta^Nk_w+\eta^Xe_w+\eta^Bb_w &&\forall w\in W,
\end{aligned}
\]

\[
\begin{aligned}
\sum_{w:(o,w,p)\in A^{OD}}x^{OD}_{owpt}
 +\sum_{c:(o,c,p)\in A^{OC}}x^{OC}_{ocpt}
 &=a_{opt} &&\text{(D4)}\\
 &\quad\forall(o,p,t)\in O\times P\times T,\\
B^{in}_{wpt}&=\sum_{o:(o,w,p)\in A^{OD}}x^{OD}_{owpt}
 +\sum_{v:(v,w,p)\in A^{DD}}x^{DD}_{vwpt}
 &&\text{(D5)}\\
 &\quad\forall(w,p,t)\in W\times P\times T,\\
B^{out}_{wpt}&=\sum_{c:(w,c,p)\in A^{DC}}x^{DC}_{wcpt}
 +\sum_{v:(w,v,p)\in A^{DD}}x^{DD}_{wvpt}\\
 &\quad\forall(w,p,t)\in W\times P\times T,\\
I_{wpt}&=I_{wp,t-1}+B^{in}_{wpt}-B^{out}_{wpt},\quad I_{wp,0}=I^0_{wp}
 &&\text{(D6)}\\
 &\quad\forall(w,p,t)\in W\times P\times T.
\end{aligned}
\]

\[
\begin{aligned}
\sum_{w:(w,c,p)\in A^{DC}}x^{DC}_{wcpt}
 +\sum_{o:(o,c,p)\in A^{OC}}x^{OC}_{ocpt}+u_{cpt}
 &=D_{cpt} &&\text{(D7)}\\
 &\quad\forall(c,p,t)\in C^D\times P\times T,\\
\sum_{w:(w,c,p)\in A^{DC}}x^{DC}_{wcpt}
 +\sum_{o:(o,c,p)\in A^{OC}}x^{OC}_{ocpt}
 &\le E_{cpt} &&\text{(D8)}\\
 &\quad\forall(c,p,t)\in C^X\times P\times T.
\end{aligned}
\]

\[
\begin{aligned}
\sum_{p\in P}I_{wpt}&\le K_w(X)+q^K_{wt} &&\forall(w,t)\in W\times T, &&\text{(D9)}\\
\sum_{p\in P}B^{in}_{wpt}&\le d_tR_w(X)+q^R_{wt} &&\forall(w,t)\in W\times T,\\
\sum_{p\in P}B^{out}_{wpt}&\le d_tH_w(X) &&\forall(w,t)\in W\times T,\\
y_w=0&\Longrightarrow q^K_{wt}=q^R_{wt}=0
 &&\forall(w,t)\in W^N\times T, &&\text{(D10)}\\
y_w,z_w,g_w&\in\{0,1\};\quad k_w,e_w,b_w,x,I,u,q^K,q^R\ge0 &&\text{(D11)}
\end{aligned}
\]

The domains in D11 are specified as follows: \(y_w\in\{0,1\}\) for all
\(w\in W^N\), \(z_w\in\{0,1\}\) for all \(w\in W^X\) and
\(g_w\in\{0,1\}\) for all \(w\in W^B\). Nonnegative \(k_w,e_w,b_w\) use
those same respective site sets. For every arc family
\(F\in\{OD,DC,DD,OC\}\), \(x^F_{ijpt}\ge0\) for all
\((i,j,p)\in A^F,t\in T\). Further, \(I_{wpt}\ge0\) for all
\((w,p,t)\in W\times P\times T\), \(u_{cpt}\ge0\) for all
\((c,p,t)\in C^D\times P\times T\), and \(q^K_{wt},q^R_{wt}\ge0\) for
all \((w,t)\in W\times T\).

D1 charges the shared investment once and explicitly includes every
transport, handling, storage and penalty term. D2 couples investment
quantities to eligible activations and makes expansion/conversion
mutually exclusive at common eligible sites. D3 distinguishes stock and
daily handling capacity. D4 allocates all exogenous supply. D5/D6
conserve grain locally with its actual initial stock. D7 permits
domestic shortage, and D8 bounds optional exports. An infinite export
upper bound is interpreted as unrestricted, consistently with the
selected input. D9 shares capacity across products: receiving overflow
is added after conversion to period tonnes. Shipping has no slack. D10
preserves the native activation indicators. D11 completes the domains.
No zero terminal-stock constraint or new minimum positive investment is
introduced.

\subsection{Two-stage stochastic model
(STO)}\label{two-stage-stochastic-model-sto}

Choose \(X=(y,z,g,k,e,b)\) once before uncertainty, imposing D2/D3 and
its domains once. For each \(s\in S\), instantiate
\(Y^s=(x^s,I^s,u^s,q^{K,s},q^{R,s})\). All routes, investment costs,
operating-day factors and penalty coefficients stay fixed. The
first-stage vector is common to every scenario.

\[
\begin{aligned}
\min_{X,\{Y^s\}_{s\in S}} Z_{STO}={}&C^{INV}(X)\\
&+\sum_{s\in S}\pi_s
 [C^{TR}(Y^s)+C^{HOLD}(Y^s)+C^{PEN}(Y^s)]. &&\text{(S1)}
\end{aligned}
\]

\[
\begin{aligned}
\sum_{w:(o,w,p)\in A^{OD}}x^{OD,s}_{owpt}
 +\sum_{c:(o,c,p)\in A^{OC}}x^{OC,s}_{ocpt}
 &=a^s_{opt} &&\text{(S2)}\\
 &\quad\forall(o,p,t,s)\in O\times P\times T\times S,\\
B^{in,s}_{wpt}&=\sum_{o:(o,w,p)\in A^{OD}}x^{OD,s}_{owpt}
 +\sum_{v:(v,w,p)\in A^{DD}}x^{DD,s}_{vwpt}
 &&\text{(S3)}\\
 &\quad\forall(w,p,t,s)\in W\times P\times T\times S,\\
B^{out,s}_{wpt}&=\sum_{c:(w,c,p)\in A^{DC}}x^{DC,s}_{wcpt}
 +\sum_{v:(w,v,p)\in A^{DD}}x^{DD,s}_{wvpt}\\
 &\quad\forall(w,p,t,s)\in W\times P\times T\times S,\\
I^s_{wpt}&=I^s_{wp,t-1}+B^{in,s}_{wpt}-B^{out,s}_{wpt},\quad I^s_{wp,0}=I^0_{wp}
 &&\text{(S4)}\\
 &\quad\forall(w,p,t,s)\in W\times P\times T\times S.
\end{aligned}
\]

\[
\begin{aligned}
\sum_{w:(w,c,p)\in A^{DC}}x^{DC,s}_{wcpt}
 +\sum_{o:(o,c,p)\in A^{OC}}x^{OC,s}_{ocpt}+u^s_{cpt}
 &=D^s_{cpt} &&\text{(S5)}\\
 &\quad\forall(c,p,t,s)\in C^D\times P\times T\times S,\\
\sum_{w:(w,c,p)\in A^{DC}}x^{DC,s}_{wcpt}
 +\sum_{o:(o,c,p)\in A^{OC}}x^{OC,s}_{ocpt}
 &\le E^s_{cpt} &&\text{(S6)}\\
 &\quad\forall(c,p,t,s)\in C^X\times P\times T\times S.
\end{aligned}
\]

\[
\begin{aligned}
\sum_{p\in P}I^s_{wpt}&\le K_w(X)+q^{K,s}_{wt}
 &&\forall(w,t,s)\in W\times T\times S, &&\text{(S7)}\\
\sum_{p\in P}B^{in,s}_{wpt}&\le d_tR_w(X)+q^{R,s}_{wt}
 &&\forall(w,t,s)\in W\times T\times S,\\
\sum_{p\in P}B^{out,s}_{wpt}&\le d_tH_w(X)
 &&\forall(w,t,s)\in W\times T\times S,\\
y_w=0&\Longrightarrow q^{K,s}_{wt}=q^{R,s}_{wt}=0
 &&\forall(w,t,s)\in W^N\times T\times S, &&\text{(S8)}\\
Y^s&\ge0 &&\forall s\in S. &&\text{(S9)}
\end{aligned}
\]

S1 applies probabilities only to recourse costs. S2--S9 are the full
scenario system; D2/D3 and common X enforce nonanticipativity. The
realised scenario gives its whole operating trajectory to recourse.
Sixty periods therefore remain a two-stage formulation, without
progressive information revelation or multistage nonanticipativity. DET
uses its original deterministic arrays with \(|S|=1\), \(\pi=1\) in this
same builder. It is not the average of the nine STO trajectories.

\subsection{Operational slacks and
penalties}\label{operational-slacks-and-penalties}

The model records unmet domestic demand, stock above static storage
capacity, and arrivals above reception capacity using separate slack
variables. They measure where the returned plan departs from the modeled
operating conditions. The objective adds these penalties to investment,
transport and inventory-holding costs. The study reports the economic subtotal
and each physical diagnostic separately.

The accepted experiments use \(\lambda^D=90000\) BRL per tonne of unmet
demand, \(\lambda^K=45000\) BRL per tonne of stock exceedance per period
and \(\lambda^R=45000\) BRL per tonne of reception overflow. Stock
exceedance accumulated over the horizon is reported in tonne-periods.
Unmet demand and excess arrivals are period flows summed in tonnes.
Appendix C describes how the implementations link these slacks to
candidate activation. Appendix D provides the unit conversions.

\subsection{Data lineage and network
construction}\label{data-lineage-and-network-construction}

The study examines one grain-storage network design problem at four
instance sizes. The instances share supply and demand data, a 60-period
planning horizon, scenario inputs, model parameters and the tables of
freight, storage and investment costs. All four retain the same 146
existing warehouses. Each larger instance adds candidate locations where
a new warehouse may be opened and retains the candidates of the smaller
instance. Table~\ref{tbl-warehouse-populations} shows how the sets
expand.

\begin{longtable}[]{@{}lrrr@{}}
\caption{Warehouse populations in the four processed
inputs.}\label{tbl-warehouse-populations}\tabularnewline
\toprule\noalign{}
Instance & Existing warehouses & Candidate locations & Total sites \\
\midrule\noalign{}
\endfirsthead
\toprule\noalign{}
Instance & Existing warehouses & Candidate locations & Total sites \\
\midrule\noalign{}
\endhead
\bottomrule\noalign{}
\endlastfoot
h215 & 146 & 69 & 215 \\
h300 & 146 & 154 & 300 \\
h400 & 146 & 254 & 400 \\
h500 & 146 & 354 & 500 \\
\end{longtable}

Each processed input workbook combines the common data with the
warehouse list and directed road distances for one instance. A facility
identifier distinguishes warehouses, including those that share
coordinates. The selection record lists the facilities in the instance
and records the stratified ranking used to add candidates. The optimizer
chooses which of the available sites to open, expand or convert to bulk
handling. Appendix D identifies the four effective workbooks, their
warehouse lists and the source data used to construct them.

The inputs include 37 distinct supply-location identifiers for soybean
and corn, 27 domestic-capital destinations and ten export-destination
roles. A city may hold both a domestic and an export role. The processed
workbooks store 63,289, 113,269, 190,569 and 287,869 directed road
distances in kilometres for h215, h300, h400 and h500, respectively.
OSRM computed the distances from a common Brazilian OpenStreetMap
extract and driving profile. The model uses a subset of routes eligible
for each product. As candidates are added, the set of admissible arcs
and the dimensions of the MILP can change. Appendix D documents the
distance calculation and arc selection.

DET uses the original supply and demand arrays. STO crosses three supply
levels (0.85, 1 and 1.15 times the base supply) with three
domestic-demand levels (0.95, 1 and 1.05 times the base demand). Each of
the nine resulting scenarios has probability 1/9. The objective measures
an undiscounted BRL total over the 60 periods.

\subsection{Experimental and validation
protocol}\label{experimental-and-validation-protocol}

For each instance, the model selects warehouse openings and expansions,
assigns capacity and plans the movement and storage of grain. The study reports
investment, transport and inventory-holding costs. Unmet demand and the two capacity exceedances are used to assess how closely the returned plan
meets the modeled operating conditions.

The h215 case provides the reference scale for comparison with the
formulation of \citet{rosa2026thesis}. It crosses DET and STO with two
network configurations. W routes grain through warehouses; D also
permits direct origin-to-customer deliveries. Two solvers are implemented in this study: a commercial solver (Gurobi\footnote{url{https://www.gurobi.com/}}), with a native implementation and a Pyomo\footnote{url{https://www.pyomo.org/}} integration; and an open-source solver (HiGHS\footnote{url{https://highs.dev/}}) only with a Pyomo implementation. The native Gurobi
comparisons use four solver threads. The h215 W STO runs through Pyomo provide a separate Gurobi-HiGHS comparison with the settings in Appendix C.

The native implementation builds and solves the model directly with gurobipy\footnote{url{https://pypi.org/project/gurobipy/}}. The native scale comparison uses W STO at h215, h300,
h400 and h500. The h300 experiment includes one observation at each of
1, 2, 4, 8 and 16 threads, and the larger native cases have one thread
each. The common-Pyomo campaign applies one experimental design to h300,
h400 and h500 W STO: two solvers, five thread ceilings and three
repetitions per combination, with seed 42, an eight-hour native time
limit and a 1\% gap target. Pyomo builds the model once per run and
passes the same algebra for a given instance through the automatic persistent Pyomo solver interface (APPSI\footnote{url{https://pyomo.readthedocs.io/en/stable/reference/topical/appsi/appsi.html}}) to Gurobi or
HiGHS. Appendix C records the settings and outcome coverage for each
experiment.

An independent check reconstructs costs and tests the exported decisions
in their original units. It uses an absolute tolerance of \(10^{-5}\)
and a relative tolerance of \(10^{-8}\) scaled to the quantities
compared. For minimization, the common relative gap is \((z-b)/|z|\) for
finite nonzero incumbent objective \(z\) and finite lower bound
\(b\le z\). The gap is zero if both values are zero. It is unavailable
for other zero incumbents, missing or nonfinite values, or inconsistent
bounds. The gap measures the quality of the objective that includes the
penalties.

\subsection{Solver parallelism and
reproducibility}\label{solver-parallelism-and-reproducibility}

The common-Pyomo campaign sets the solver thread ceiling to 1, 2, 4, 8
or 16. Gurobi uses one MIP search (\texttt{ConcurrentMIP=1}) with
parallel algorithms that can use the specified threads. HiGHS 1.15.1
uses \texttt{parallel="on"} to enable parallel tree search within its
thread ceiling. In that version, the simplex solver used for the LP
relaxations inside the MIP runs serially. The computational effect of
extra threads therefore depends on where each solver spends its time.

In branch-and-bound, the solver chooses a subproblem to explore, solves
its relaxation and uses the resulting bound and candidate solutions to
guide the next steps. The seed fixes the initialization of a solver's
pseudo-random choices. Branching, node selection, cuts and heuristics
also depend on the evolving search state. Thread settings can change
that state and the resulting trajectory. A time limit can stop repeated
runs at different points in their search. Appendix C explains the two
parallel modes, the settings used and how to interpret the repetitions.

\paragraph{Exploratory root-method experiment.} A separate CPU-only h300 W STO experiment fixes the input instance and compares four Gurobi 13.0.3 root LP methods (dual simplex, barrier, deterministic concurrent LP and PDHG), three heuristic profiles (standard, \texttt{Heuristics=0}, and standard with a fixed NoRel budget), and physical-core allocations of 1, 2, 4, 8 and 16. The first block comprises 60 planned tasks; temporal repetitions are intended but are not assumed complete. All runs retain seed 42, an eight-hour native limit, the 1\% relative-gap target and independent incumbent validation. Slurm and solver receipts distinguish terminal success, time limits, memory-guard interruption, errors and unexecuted tasks. Root-phase time, incumbent and target-attainment events, processed nodes and job-cgroup memory are interpreted only where recorded. This exploratory block is not pooled with the historical solver-thread campaign.

\subsection{HPC experimental
environment}\label{hpc-experimental-environment}

The high-performane computing (HPC) experiments use NPAD/UFRN's Intel shared-memory infrastructure under
Slurm \citep{npad2026}. Retained runtime receipts identify Intel Xeon E5-2698 v3 CPUs, Python 3.13 and Gurobi 13.0.3. The scalar allocations request 24 logical CPUs and 192 GiB. Gurobi's SoftMemLimit is 160 GB, but this is not a process-RSS ceiling. The historical native h300 worker affinity pool contains 16 logical CPU IDs across 12 physical cores, including SMT siblings. Solver threads, affinity and allocated logical CPUs are distinct registered resource scopes.

For the common-Pyomo campaign the profile requests an exclusive node and 192 GiB for a 12-hour allocation. The h300 worker affinity is 16 logical CPU IDs. Scheduler accounting records 64 allocated logical CPUs at the job level for the scientific jobs. Solver thread ceilings are settings within this envelope. The runtime and hardware receipts identify the conditions for each numerical observation.

Historical native reservations are 14 hours for the h215 pair, 30 hours for the five-arm h300 block, and eight hours each for h400 and h500.
Reservation time is not measured optimizer time. Current models solve within one allocated node, and the reported scalar blocks run serially.
Distributing independent experiments through scheduler arrays could not make an individual model a distributed MILP solve; no multi-node speedup has been claimed. Network preparation, construction and post-processing are CPU workloads outside the primary optimization call. No GPU acceleration or use of the InfiniBand fabric by a distributed MILP algorithm was
asserted.

\subsection{Model size, memory and performance
measurement}\label{model-size-memory-and-performance-measurement}

With retained product-eligible arcs \(\mathcal A_p\), flow variables
grow as \(|S||T|\sum_p|\mathcal A_p|\). Inventory contributes terms
proportional to \(|S||T||W||P|\), with separate slack and first-stage
blocks. Before screening, the directed interwarehouse graph can contain
\(|W|(|W|-1)\) arcs. A fixed retained fraction therefore does not
eliminate quadratic network growth. These are structural matrix counts,
not an asymptotic algorithmic-complexity proof.

Native runtime, model-construction wall time, application optimization
wall time, extraction, validation, export and application end-to-end
time are separate clocks. Allocation elapsed includes additional staging
and control. The microsecond data-read fields at an already-loaded model
layer are not large-workbook ETL timings. The study does not reconstruct missing
phases by subtracting overlapping clocks. Memory labels distinguish
sampled application RSS, solver-managed memory and Slurm step MaxRSS.
Sampling may miss an instantaneous peak; requested memory and allocated
CPU-hour products are not measured consumption or billing.

For a fixed instance and solver, strong scaling varies resources while
holding the problem fixed. With comparable work, timing and quality
criteria, speedup is \(S(p)=T(1)/T(p)\) and parallel efficiency is
\(E(p)=S(p)/p\) for a \(p\)-fold resource setting relative to the
baseline. In this study a ratio normalized by the configured thread
ceiling measures performance relative to that setting. Effective CPU
utilization requires a separate measurement. A time-to-target comparison
uses the recorded time when a common quality threshold is first
attained. Terminal runtime also includes work up to the observed
stopping condition.

Across populations the experiment evaluates problem-size and resource
scaling. The selected graph and the number of nonzeros change
nonproportionally, while classical weak scaling keeps the workload per
resource approximately constant. Nonzeros per thread describe structural
load. Solver parallelism concerns work within one solve. Concurrent
Slurm jobs concern independent executions. Throughput would count
specified completed work per unit elapsed time over a stated allocation
and observation window.

The initial LP-relaxation clock and branch-and-bound node count are additional phase diagnostics. A recorded one-node tree denotes the observed processed-node count under a particular algorithm and stopping rule; it is not a complexity bound on the formulation. A root-time ratio requires comparable algorithm, heuristic and resource settings and is not automatically total MILP speedup.

\subsection{Geographic figures}\label{geographic-figures}

The maps display facility positions and the investments returned by the
model on national and federative-unit boundaries from IBGE. Map lines
delineate study areas and do not necessarily depict accepted national
boundaries. Appendix E explains the symbols and links the exported flows
to their stored OSRM road distances.

\subsection{Human-supervised AI-assisted
research}\label{human-supervised-ai-assisted-research}

AI assistants supported mathematical discussion, implementation and
scientific writing under human supervision. Gurobot supported
formulation discussions; Gemini and GPT systems supported technical
criticism; ChatGPT Classic supported coordination; and Codex supported
Python development, tests and documentation. Researchers selected
assumptions and experimental configurations, reviewed changes, executed
the HPC commands and interpreted the resulting plans.

The workflow in Figure~\ref{fig-human-ai} connects human scientific
intent to implementation, verification, scientific assessment and
reporting. Appendix B explains its six phases and four complementary
verification levels.

\begin{figure}[H]

\centering{

\pandocbounded{\includegraphics[keepaspectratio]{human_ai_workflow.pdf}}

}

\caption{\label{fig-human-ai}Human-supervised research workflow.}

\end{figure}%

The return arrows represent revision when a discrepancy is found. For
example, shared investment and probability-weighted recourse were
checked in analytical examples before large experiments. Original-unit
residuals and objective reconstruction then supported assessment of the
returned decisions. Each result was associated with its source, inputs,
solver configuration and validation record. Human scientific judgment
connected this evidence to conclusions about computational requirements
and agricultural planning.

\section{Results}\label{results}

The h215 comparison examines the effect of direct deliveries and
scenario-dependent operation. The historical native W STO series then
traces how investment choices and computational effort change as
candidate sites are added. The repeated common-Pyomo experiments provide
a fixed-instance comparison of thread ceilings at h300, followed by a
complete single-repetition thread grid at h400.

\subsection{h215 comparison of warehouse-only and direct-enabled
networks}\label{h215-comparison-of-warehouse-only-and-direct-enabled-networks}

The h215 comparison uses one set of sites, input data and costs in all
four cases. D adds direct origin-to-customer routes to the warehouse
routes shared with W. The native gurobipy runs used four solver threads.
Appendix C gives their settings.

The h215 results show how planning under uncertainty and allowing direct
deliveries change warehouse investments and economic costs. The W STO
series then shows how the plans change as the same network design
problem is evaluated with more candidate locations. The solver
comparisons use the Pyomo runs and are reported separately.
Table~\ref{tbl-wd-h215} compares the h215 warehouse-only and
direct-enabled plans under deterministic and stochastic planning.
Optimal native status means the configured gap criterion, not exact
zero-gap optimality.

\begin{longtable}[]{@{}
  >{\raggedright\arraybackslash}p{(\linewidth - 10\tabcolsep) * \real{0.1667}}
  >{\raggedleft\arraybackslash}p{(\linewidth - 10\tabcolsep) * \real{0.1667}}
  >{\raggedleft\arraybackslash}p{(\linewidth - 10\tabcolsep) * \real{0.1667}}
  >{\raggedleft\arraybackslash}p{(\linewidth - 10\tabcolsep) * \real{0.1667}}
  >{\raggedleft\arraybackslash}p{(\linewidth - 10\tabcolsep) * \real{0.1667}}
  >{\raggedleft\arraybackslash}p{(\linewidth - 10\tabcolsep) * \real{0.1667}}@{}}
\caption{Original validated h215 W/D
observations.}\label{tbl-wd-h215}\tabularnewline
\toprule\noalign{}
\begin{minipage}[b]{\linewidth}\raggedright
Mode/\allowbreak topology
\end{minipage} & \begin{minipage}[b]{\linewidth}\raggedleft
Objective, BRL trillion
\end{minipage} & \begin{minipage}[b]{\linewidth}\raggedleft
Economic subtotal, BRL billion
\end{minipage} & \begin{minipage}[b]{\linewidth}\raggedleft
Gap, \%
\end{minipage} & \begin{minipage}[b]{\linewidth}\raggedleft
Native Runtime, min
\end{minipage} & \begin{minipage}[b]{\linewidth}\raggedleft
Domestic service, \%
\end{minipage} \\
\midrule\noalign{}
\endfirsthead
\toprule\noalign{}
\begin{minipage}[b]{\linewidth}\raggedright
Mode/\allowbreak topology
\end{minipage} & \begin{minipage}[b]{\linewidth}\raggedleft
Objective, BRL trillion
\end{minipage} & \begin{minipage}[b]{\linewidth}\raggedleft
Economic subtotal, BRL billion
\end{minipage} & \begin{minipage}[b]{\linewidth}\raggedleft
Gap, \%
\end{minipage} & \begin{minipage}[b]{\linewidth}\raggedleft
Native Runtime, min
\end{minipage} & \begin{minipage}[b]{\linewidth}\raggedleft
Domestic service, \%
\end{minipage} \\
\midrule\noalign{}
\endhead
\bottomrule\noalign{}
\endlastfoot
DET W & 11.437342 & 389.810 & 0.000812 & 2.04 & 91.861 \\
DET D & 7.340492 & 338.715 & 0.552427 & 1.73 & 89.896 \\
STO W & 13.071397 & 393.902 & 0.740159 & 40.76 & 92.267 \\
STO D & 7.355002 & 340.109 & 0.623718 & 43.04 & 89.857 \\
\end{longtable}

Note that the economic subtotal is only in the order of 3\% for W and
5\% for D of the objective function value, due to the significant
penalties included. The direct-enabled plans reduce economic costs by
13.11\% in DET and 13.66\% in STO. All 69 candidates receive positive
capacity in W and D. With direct deliveries, 142 existing sites are
expanded in each plan, compared with 80 in DET W and 75 in STO W. Four
existing sites are converted for bulk handling in each D plan, compared
with 66 and 71 in the corresponding W plans.

Direct routes carry 82.17\% of destination arrivals in DET D and 81.81\%
in STO D. Domestic service declines from 91.861\% to 89.896\% in DET and
from 92.267\% to 89.857\% in STO. The D plans therefore combine lower
economic costs and more static expansion with greater unmet domestic
demand. Appendix C gives the associated bounds and effort measurements.

\subsection{Investment and operational interpretation across
populations}\label{investment-and-operational-interpretation-across-populations}

Table~\ref{tbl-scalar-financial-components} separates investment and
operating costs. Investment is charged once; STO operation is
scenario-weighted over the horizon. No annualization, currency
conversion or price-year adjustment is applied.

\begin{longtable}[]{@{}
  >{\raggedright\arraybackslash}p{(\linewidth - 6\tabcolsep) * \real{0.2500}}
  >{\raggedleft\arraybackslash}p{(\linewidth - 6\tabcolsep) * \real{0.2500}}
  >{\raggedleft\arraybackslash}p{(\linewidth - 6\tabcolsep) * \real{0.2500}}
  >{\raggedleft\arraybackslash}p{(\linewidth - 6\tabcolsep) * \real{0.2500}}@{}}
\caption{Reconstructed financial components, excluding
penalties.}\label{tbl-scalar-financial-components}\tabularnewline
\toprule\noalign{}
\begin{minipage}[b]{\linewidth}\raggedright
Instance
\end{minipage} & \begin{minipage}[b]{\linewidth}\raggedleft
One-time investment, BRL billion
\end{minipage} & \begin{minipage}[b]{\linewidth}\raggedleft
Transport, BRL billion
\end{minipage} & \begin{minipage}[b]{\linewidth}\raggedleft
Inventory holding, BRL billion
\end{minipage} \\
\midrule\noalign{}
\endfirsthead
\toprule\noalign{}
\begin{minipage}[b]{\linewidth}\raggedright
Instance
\end{minipage} & \begin{minipage}[b]{\linewidth}\raggedleft
One-time investment, BRL billion
\end{minipage} & \begin{minipage}[b]{\linewidth}\raggedleft
Transport, BRL billion
\end{minipage} & \begin{minipage}[b]{\linewidth}\raggedleft
Inventory holding, BRL billion
\end{minipage} \\
\midrule\noalign{}
\endhead
\bottomrule\noalign{}
\endlastfoot
h215 DET & 9.848 & 364.960 & 15.002 \\
h215 STO & 9.819 & 367.138 & 16.946 \\
h300 STO & 11.260 & 360.578 & 13.664 \\
h400 STO & 12.539 & 356.123 & 12.911 \\
h500 STO & 14.318 & 346.074 & 12.781 \\
\end{longtable}

Table~\ref{tbl-scalar-investments} and Table~\ref{tbl-scalar-quantities}
report the locations and quantities of investments. Counts are obtained
from the accepted original decision vectors, using a display threshold
of 1e-8 t for h215--h400 and 1e-7 t for the original h500 aggregate.
Existing availability is not counted as candidate opening.

\begin{longtable}[]{@{}
  >{\raggedright\arraybackslash}p{(\linewidth - 8\tabcolsep) * \real{0.2000}}
  >{\raggedleft\arraybackslash}p{(\linewidth - 8\tabcolsep) * \real{0.2000}}
  >{\raggedleft\arraybackslash}p{(\linewidth - 8\tabcolsep) * \real{0.2000}}
  >{\raggedleft\arraybackslash}p{(\linewidth - 8\tabcolsep) * \real{0.2000}}
  >{\raggedleft\arraybackslash}p{(\linewidth - 8\tabcolsep) * \real{0.2000}}@{}}
\caption{Counts from the accepted original decision
vectors.}\label{tbl-scalar-investments}\tabularnewline
\toprule\noalign{}
\begin{minipage}[b]{\linewidth}\raggedright
Instance
\end{minipage} & \begin{minipage}[b]{\linewidth}\raggedleft
Available candidates
\end{minipage} & \begin{minipage}[b]{\linewidth}\raggedleft
Candidates with positive capacity
\end{minipage} & \begin{minipage}[b]{\linewidth}\raggedleft
Expanded existing sites
\end{minipage} & \begin{minipage}[b]{\linewidth}\raggedleft
Bulkified existing sites
\end{minipage} \\
\midrule\noalign{}
\endfirsthead
\toprule\noalign{}
\begin{minipage}[b]{\linewidth}\raggedright
Instance
\end{minipage} & \begin{minipage}[b]{\linewidth}\raggedleft
Available candidates
\end{minipage} & \begin{minipage}[b]{\linewidth}\raggedleft
Candidates with positive capacity
\end{minipage} & \begin{minipage}[b]{\linewidth}\raggedleft
Expanded existing sites
\end{minipage} & \begin{minipage}[b]{\linewidth}\raggedleft
Bulkified existing sites
\end{minipage} \\
\midrule\noalign{}
\endhead
\bottomrule\noalign{}
\endlastfoot
h215 DET & 69 & 69 & 80 & 66 \\
h215 STO & 69 & 69 & 75 & 71 \\
h300 STO & 154 & 154 & 70 & 76 \\
h400 STO & 254 & 254 & 69 & 77 \\
h500 STO & 354 & 354 & 72 & 74 \\
\end{longtable}

The bulkification basis determines handling-rate gains; it is not
additional static storage.

\begin{longtable}[]{@{}
  >{\raggedright\arraybackslash}p{(\linewidth - 6\tabcolsep) * \real{0.2500}}
  >{\raggedleft\arraybackslash}p{(\linewidth - 6\tabcolsep) * \real{0.2500}}
  >{\raggedleft\arraybackslash}p{(\linewidth - 6\tabcolsep) * \real{0.2500}}
  >{\raggedleft\arraybackslash}p{(\linewidth - 6\tabcolsep) * \real{0.2500}}@{}}
\caption{Investment
quantities.}\label{tbl-scalar-quantities}\tabularnewline
\toprule\noalign{}
\begin{minipage}[b]{\linewidth}\raggedright
Instance
\end{minipage} & \begin{minipage}[b]{\linewidth}\raggedleft
Candidate capacity, million t
\end{minipage} & \begin{minipage}[b]{\linewidth}\raggedleft
Static expansion, million t
\end{minipage} & \begin{minipage}[b]{\linewidth}\raggedleft
Bulkification investment basis, million t
\end{minipage} \\
\midrule\noalign{}
\endfirsthead
\toprule\noalign{}
\begin{minipage}[b]{\linewidth}\raggedright
Instance
\end{minipage} & \begin{minipage}[b]{\linewidth}\raggedleft
Candidate capacity, million t
\end{minipage} & \begin{minipage}[b]{\linewidth}\raggedleft
Static expansion, million t
\end{minipage} & \begin{minipage}[b]{\linewidth}\raggedleft
Bulkification investment basis, million t
\end{minipage} \\
\midrule\noalign{}
\endhead
\bottomrule\noalign{}
\endlastfoot
h215 DET & 7.408444 & 0.395351 & 0.445250 \\
h215 STO & 7.381322 & 0.391226 & 0.461750 \\
h300 STO & 8.522815 & 0.379411 & 0.509010 \\
h400 STO & 9.552150 & 0.385596 & 0.484270 \\
h500 STO & 10.992729 & 0.389176 & 0.437018 \\
\end{longtable}

Table~\ref{tbl-scalar-physical} gives the corresponding service, stock
and reception diagnostics. STO quantities are probability-weighted, not
unweighted sums over nine scenarios.

\begin{longtable}[]{@{}
  >{\raggedright\arraybackslash}p{(\linewidth - 8\tabcolsep) * \real{0.2000}}
  >{\raggedleft\arraybackslash}p{(\linewidth - 8\tabcolsep) * \real{0.2000}}
  >{\raggedleft\arraybackslash}p{(\linewidth - 8\tabcolsep) * \real{0.2000}}
  >{\raggedleft\arraybackslash}p{(\linewidth - 8\tabcolsep) * \real{0.2000}}
  >{\raggedleft\arraybackslash}p{(\linewidth - 8\tabcolsep) * \real{0.2000}}@{}}
\caption{Physical diagnostics over the declared
horizon.}\label{tbl-scalar-physical}\tabularnewline
\toprule\noalign{}
\begin{minipage}[b]{\linewidth}\raggedright
Instance
\end{minipage} & \begin{minipage}[b]{\linewidth}\raggedleft
Domestic service, \%
\end{minipage} & \begin{minipage}[b]{\linewidth}\raggedleft
Unmet demand, million t
\end{minipage} & \begin{minipage}[b]{\linewidth}\raggedleft
Stock exceedance, million t period
\end{minipage} & \begin{minipage}[b]{\linewidth}\raggedleft
Reception overflow, million t
\end{minipage} \\
\midrule\noalign{}
\endfirsthead
\toprule\noalign{}
\begin{minipage}[b]{\linewidth}\raggedright
Instance
\end{minipage} & \begin{minipage}[b]{\linewidth}\raggedleft
Domestic service, \%
\end{minipage} & \begin{minipage}[b]{\linewidth}\raggedleft
Unmet demand, million t
\end{minipage} & \begin{minipage}[b]{\linewidth}\raggedleft
Stock exceedance, million t period
\end{minipage} & \begin{minipage}[b]{\linewidth}\raggedleft
Reception overflow, million t
\end{minipage} \\
\midrule\noalign{}
\endhead
\bottomrule\noalign{}
\endlastfoot
h215 DET & 91.8612 & 55.1578 & 77.8448 & 57.3404 \\
h215 STO & 92.2674 & 52.4049 & 109.0691 & 67.8432 \\
h300 STO & 91.3256 & 58.7873 & 46.4440 & 33.8392 \\
h400 STO & 91.3304 & 58.7551 & 26.2548 & 16.4977 \\
h500 STO & 91.6360 & 56.6840 & 14.2025 & 4.0754 \\
\end{longtable}

All included candidates are activated and receive positive capacity
through h500. The h500 aggregate records 354 candidate investments, 72
positive expansions and 74 positive bulk conversions. The original
h215--h400 vector audits additionally establish positive transport at
every included warehouse and nine used export roles; those two
geographic-use findings are not extrapolated to the h500 aggregate.
These are returned plan characteristics, not uniqueness or
infrastructure feasibility.

Across the returned STO plans, additional candidate sites accompany more
installed candidate capacity and greater one-time investment, while the
transport and holding components are lower. For example, from h300 to
h500, investment increases from BRL 11.26 to 14.32 billion and transport
costs fall from BRL 360.58 to 346.07 billion in the returned plans. The
physical diagnostics show how those plans balance unmet demand and
capacity exceedances.

Besides, across the W STO plans, positive candidate capacity rises from
7.38 million tonnes at h215 to 10.99 million tonnes at h500. Static
expansion stays near 0.38--0.39 million tonnes. The added storage is
concentrated in new candidate warehouses, while handling capacity and
domestic service describe separate aspects of the plans.

\subsubsection{Geographic
interpretation}\label{geographic-interpretation}

Figure~\ref{fig-map-h215}, Figure~\ref{fig-map-h300},
Figure~\ref{fig-map-h400} and Figure~\ref{fig-map-h500} locate the
investments in the W STO plans. The maps share the Albers equal-area
conic projection and retain original facility coordinates. The inputs
share the model and economic parameters while expanding the set of
candidates. The maps show investment locations; route lines are omitted
for legibility. Appendix D identifies the stored OSRM distances, and
Appendix E links them to the positive flows in the returned plans.

\begin{figure}[h]

\centering{

\pandocbounded{\includegraphics[keepaspectratio]{scalar_h215_sto.pdf}}

}

\caption{\label{fig-map-h215}Stochastic plan at 215 warehouse sites.}

\end{figure}%

\begin{figure}[H]

\centering{

\pandocbounded{\includegraphics[keepaspectratio]{scalar_h300_t1.pdf}}

}

\caption{\label{fig-map-h300}Stochastic plan at 300 warehouse sites.}

\end{figure}%

\begin{figure}[h]

\centering{

\pandocbounded{\includegraphics[keepaspectratio]{scalar_h400_t1.pdf}}

}

\caption{\label{fig-map-h400}Stochastic plan at 400 warehouse sites.}

\end{figure}%

\begin{figure}[H]

\centering{

\pandocbounded{\includegraphics[keepaspectratio]{scalar_h500_t1.pdf}}

}

\caption{\label{fig-map-h500}Stochastic investments at 500 warehouse
sites.}

\end{figure}%

The original h500 decision vector assigns positive capacity to 354
candidates, expands 72 existing sites and converts 74 for bulk handling.
Appendix E provides the capacity maps and distinguishes these investment
decisions from the operating flows.

\subsection{Native Gurobi computation across warehouse
populations}\label{native-gurobi-computation-across-warehouse-populations}

Table~\ref{tbl-scalar-size}, Table~\ref{tbl-scalar-resources} and
Table~\ref{tbl-scalar-quality} describe the computational size, effort
and solution quality of the native runs.

The original h500 result extends the W STO series to 500 sites. The
representative native results use the validated h215 DET/STO
observations and the one-thread W STO observations at h300, h400 and
h500.

\begin{longtable}[]{@{}lrrrr@{}}
\caption{Original model
dimensions.}\label{tbl-scalar-size}\tabularnewline
\toprule\noalign{}
Instance & Variables & Binary & Linear constraints & Matrix nonzeros \\
\midrule\noalign{}
\endfirsthead
\toprule\noalign{}
Instance & Variables & Binary & Linear constraints & Matrix nonzeros \\
\midrule\noalign{}
\endhead
\bottomrule\noalign{}
\endlastfoot
h215 DET & 1,562,174 & 307 & 74,017 & 5,786,346 \\
h215 STO & 14,054,654 & 307 & 661,057 & 52,068,826 \\
h300 STO & 25,107,544 & 392 & 890,642 & 94,635,126 \\
h400 STO & 42,106,944 & 492 & 1,160,742 & 160,697,126 \\
h500 STO & 63,426,344 & 592 & 1,430,842 & 244,039,126 \\
\end{longtable}

General indicator constraints are additional and remain separately
recorded; binary counts are included in integer counts.

\begin{longtable}[]{@{}
  >{\raggedright\arraybackslash}p{(\linewidth - 8\tabcolsep) * \real{0.2000}}
  >{\raggedleft\arraybackslash}p{(\linewidth - 8\tabcolsep) * \real{0.2000}}
  >{\raggedleft\arraybackslash}p{(\linewidth - 8\tabcolsep) * \real{0.2000}}
  >{\raggedleft\arraybackslash}p{(\linewidth - 8\tabcolsep) * \real{0.2000}}
  >{\raggedleft\arraybackslash}p{(\linewidth - 8\tabcolsep) * \real{0.2000}}@{}}
\caption{Measured phases and memory
scopes.}\label{tbl-scalar-resources}\tabularnewline
\toprule\noalign{}
\begin{minipage}[b]{\linewidth}\raggedright
Instance
\end{minipage} & \begin{minipage}[b]{\linewidth}\raggedleft
Native Runtime, min
\end{minipage} & \begin{minipage}[b]{\linewidth}\raggedleft
Build, min
\end{minipage} & \begin{minipage}[b]{\linewidth}\raggedleft
Application end-to-end, min
\end{minipage} & \begin{minipage}[b]{\linewidth}\raggedleft
Sampled application RSS, GiB
\end{minipage} \\
\midrule\noalign{}
\endfirsthead
\toprule\noalign{}
\begin{minipage}[b]{\linewidth}\raggedright
Instance
\end{minipage} & \begin{minipage}[b]{\linewidth}\raggedleft
Native Runtime, min
\end{minipage} & \begin{minipage}[b]{\linewidth}\raggedleft
Build, min
\end{minipage} & \begin{minipage}[b]{\linewidth}\raggedleft
Application end-to-end, min
\end{minipage} & \begin{minipage}[b]{\linewidth}\raggedleft
Sampled application RSS, GiB
\end{minipage} \\
\midrule\noalign{}
\endhead
\bottomrule\noalign{}
\endlastfoot
h215 DET & 2.04 & 0.41 & 2.91 & 5.54 \\
h215 STO & 40.76 & 3.98 & 46.22 & 43.13 \\
h300 STO & 47.14 & 7.19 & 57.09 & 46.35 \\
h400 STO & 70.46 & 11.87 & 87.02 & 72.34 \\
h500 STO & 127.58 & 18.16 & 152.72 & 111.78 \\
\end{longtable}

Allocation elapsed, Slurm MaxRSS and solver-managed memory are different
measurements.

\begin{longtable}[]{@{}
  >{\raggedright\arraybackslash}p{(\linewidth - 8\tabcolsep) * \real{0.2000}}
  >{\raggedleft\arraybackslash}p{(\linewidth - 8\tabcolsep) * \real{0.2000}}
  >{\raggedleft\arraybackslash}p{(\linewidth - 8\tabcolsep) * \real{0.2000}}
  >{\raggedleft\arraybackslash}p{(\linewidth - 8\tabcolsep) * \real{0.2000}}
  >{\raggedleft\arraybackslash}p{(\linewidth - 8\tabcolsep) * \real{0.2000}}@{}}
\caption{Scalar solution quality and cost
composition.}\label{tbl-scalar-quality}\tabularnewline
\toprule\noalign{}
\begin{minipage}[b]{\linewidth}\raggedright
Instance
\end{minipage} & \begin{minipage}[b]{\linewidth}\raggedleft
Penalized objective, BRL trillion
\end{minipage} & \begin{minipage}[b]{\linewidth}\raggedleft
Economic subtotal, BRL billion
\end{minipage} & \begin{minipage}[b]{\linewidth}\raggedleft
Artificial penalties, \%
\end{minipage} & \begin{minipage}[b]{\linewidth}\raggedleft
Global gap, \%
\end{minipage} \\
\midrule\noalign{}
\endfirsthead
\toprule\noalign{}
\begin{minipage}[b]{\linewidth}\raggedright
Instance
\end{minipage} & \begin{minipage}[b]{\linewidth}\raggedleft
Penalized objective, BRL trillion
\end{minipage} & \begin{minipage}[b]{\linewidth}\raggedleft
Economic subtotal, BRL billion
\end{minipage} & \begin{minipage}[b]{\linewidth}\raggedleft
Artificial penalties, \%
\end{minipage} & \begin{minipage}[b]{\linewidth}\raggedleft
Global gap, \%
\end{minipage} \\
\midrule\noalign{}
\endhead
\bottomrule\noalign{}
\endlastfoot
h215 DET & 11.437342 & 389.810 & 96.592 & 0.000812 \\
h215 STO & 13.071397 & 393.902 & 96.987 & 0.740159 \\
h300 STO & 9.289109 & 385.502 & 95.850 & 0.004119 \\
h400 STO & 7.593395 & 381.573 & 94.975 & 0.836784 \\
h500 STO & 6.297240 & 373.173 & 94.074 & 0.207432 \\
\end{longtable}

Native optimal status denotes the configured gap criterion, not exact
zero-gap proof.

\subsection{Descriptive thread observations at
h300}\label{descriptive-thread-observations-at-h300}

Table~\ref{tbl-scalar-threads} and Figure~\ref{fig-scalar-threads}
summarize the five native h300 settings by solver thread count.

\begin{longtable}[]{@{}
  >{\raggedleft\arraybackslash}p{(\linewidth - 6\tabcolsep) * \real{0.2500}}
  >{\raggedleft\arraybackslash}p{(\linewidth - 6\tabcolsep) * \real{0.2500}}
  >{\raggedleft\arraybackslash}p{(\linewidth - 6\tabcolsep) * \real{0.2500}}
  >{\raggedleft\arraybackslash}p{(\linewidth - 6\tabcolsep) * \real{0.2500}}@{}}
\caption{Native gurobipy results for h300 W STO, one observation per
thread setting.}\label{tbl-scalar-threads}\tabularnewline
\toprule\noalign{}
\begin{minipage}[b]{\linewidth}\raggedleft
Solver threads
\end{minipage} & \begin{minipage}[b]{\linewidth}\raggedleft
Native Runtime, min
\end{minipage} & \begin{minipage}[b]{\linewidth}\raggedleft
Sampled application RSS, GiB
\end{minipage} & \begin{minipage}[b]{\linewidth}\raggedleft
Global gap, \%
\end{minipage} \\
\midrule\noalign{}
\endfirsthead
\toprule\noalign{}
\begin{minipage}[b]{\linewidth}\raggedleft
Solver threads
\end{minipage} & \begin{minipage}[b]{\linewidth}\raggedleft
Native Runtime, min
\end{minipage} & \begin{minipage}[b]{\linewidth}\raggedleft
Sampled application RSS, GiB
\end{minipage} & \begin{minipage}[b]{\linewidth}\raggedleft
Global gap, \%
\end{minipage} \\
\midrule\noalign{}
\endhead
\bottomrule\noalign{}
\endlastfoot
1 & 47.143 & 46.352 & 0.004119 \\
2 & 53.008 & 57.244 & 0.004119 \\
4 & 58.888 & 77.958 & 0.004119 \\
8 & 67.454 & 78.211 & 0.004119 \\
16 & 67.613 & 78.368 & 0.004119 \\
\end{longtable}

All five h300 objectives and best bounds coincide; their original
warehouse, flow, inventory and shortage CSVs are byte-identical across
arms. Native Runtime ranges from 47.14 to 67.61 minutes and sampled
application RSS from 46.35 to 78.37 GiB. The single-thread run had the
lowest recorded time and sampled memory. Figure~\ref{fig-scalar-threads}
compares the observed runtime and sampled memory for these five
configurations.

\begin{figure}[h]

\centering{

\pandocbounded{\includegraphics[keepaspectratio]{scalar_threads.pdf}}

}

\caption{\label{fig-scalar-threads}Thread-count observations at 300
sites.}

\end{figure}%

\subsection{h215 W STO with the common Pyomo
implementation}\label{h215-w-sto-with-the-common-pyomo-implementation}

For the h215 Pyomo runs, Table~\ref{tbl-pyomo-settings} gives the
settings, Table~\ref{tbl-highs-h215} gives solution quality and native
time, and Table~\ref{tbl-pyomo-outcomes} gives application time and
sampled memory.

The completed h215 W STO calculations use the same frozen workbook,
ordered warehouses, selected directed network, nine equiprobable
scenarios and scalar coefficients. Both Pyomo pipelines use the bounded
candidate-slack representation of Appendix C. This gives a more specific
comparison basis than the h215 label alone. The earlier HiGHS record has
no ordered-algebra fingerprint; its retained input bindings, model
dimensions and bound policy agree with the common formulation, but an
unrecorded fingerprint is not inferred retrospectively.

\begin{longtable}[]{@{}
  >{\raggedright\arraybackslash}p{(\linewidth - 4\tabcolsep) * \real{0.3333}}
  >{\raggedright\arraybackslash}p{(\linewidth - 4\tabcolsep) * \real{0.3333}}
  >{\raggedright\arraybackslash}p{(\linewidth - 4\tabcolsep) * \real{0.3333}}@{}}
\caption{Effective configurations of the completed Pyomo h215 W STO
observations.}\label{tbl-pyomo-settings}\tabularnewline
\toprule\noalign{}
\begin{minipage}[b]{\linewidth}\raggedright
Setting
\end{minipage} & \begin{minipage}[b]{\linewidth}\raggedright
Gurobi / APPSI
\end{minipage} & \begin{minipage}[b]{\linewidth}\raggedright
HiGHS / APPSI
\end{minipage} \\
\midrule\noalign{}
\endfirsthead
\toprule\noalign{}
\begin{minipage}[b]{\linewidth}\raggedright
Setting
\end{minipage} & \begin{minipage}[b]{\linewidth}\raggedright
Gurobi / APPSI
\end{minipage} & \begin{minipage}[b]{\linewidth}\raggedright
HiGHS / APPSI
\end{minipage} \\
\midrule\noalign{}
\endhead
\bottomrule\noalign{}
\endlastfoot
Native version & 13.0.3 & 1.15.1 \\
Pyomo version & 6.10.1 & 6.10.1 \\
Native time allowance, s & 28,800 & 28,800 \\
Relative MIP gap target & 0.01 & 0.01 \\
Thread ceiling & 4 & 4 \\
Seed setting & Seed = 42 & random\_seed = 42 \\
Feasibility controls & FeasibilityTol and IntFeasTol = 1e-9 &
mip\_feasibility\_tolerance = 1e-9 \\
\end{longtable}

The historical HiGHS log records parallel tree search off; this differs
from the h300--h500 parallel-search profile. The Gurobi calculation
supplied an ordered pre-solve algebra fingerprint for 14,054,654
variables, 307 binaries, 735,577 linear constraints and 52,217,866
nonzeros. The earlier HiGHS model has the same dimensions. Matching
dimensions alone do not prove identical coefficients or identical native
presolved models. Original input and mathematical-contract checks
therefore accompany the dimensional comparison. Both returned incumbents
passed their original independent primal validation and the prescribed
1\% scalar-gap criterion.

\begin{longtable}[]{@{}
  >{\raggedright\arraybackslash}p{(\linewidth - 6\tabcolsep) * \real{0.2500}}
  >{\raggedleft\arraybackslash}p{(\linewidth - 6\tabcolsep) * \real{0.2500}}
  >{\raggedleft\arraybackslash}p{(\linewidth - 6\tabcolsep) * \real{0.2500}}
  >{\raggedleft\arraybackslash}p{(\linewidth - 6\tabcolsep) * \real{0.2500}}@{}}
\caption{Attained scalar quality and native computation time for the two
completed Pyomo observations.}\label{tbl-highs-h215}\tabularnewline
\toprule\noalign{}
\begin{minipage}[b]{\linewidth}\raggedright
Implementation
\end{minipage} & \begin{minipage}[b]{\linewidth}\raggedleft
Objective, BRL trillion
\end{minipage} & \begin{minipage}[b]{\linewidth}\raggedleft
Common gap, \%
\end{minipage} & \begin{minipage}[b]{\linewidth}\raggedleft
Native runtime, min
\end{minipage} \\
\midrule\noalign{}
\endfirsthead
\toprule\noalign{}
\begin{minipage}[b]{\linewidth}\raggedright
Implementation
\end{minipage} & \begin{minipage}[b]{\linewidth}\raggedleft
Objective, BRL trillion
\end{minipage} & \begin{minipage}[b]{\linewidth}\raggedleft
Common gap, \%
\end{minipage} & \begin{minipage}[b]{\linewidth}\raggedleft
Native runtime, min
\end{minipage} \\
\midrule\noalign{}
\endhead
\bottomrule\noalign{}
\endlastfoot
Pyomo/APPSI Gurobi & 13.071397 & 0.740159 & 31.54 \\
Pyomo/APPSI HiGHS & 12.974673 & 0.000196 & 426.70 \\
\end{longtable}

These are separate single runs, not matched repeated measurements at
equal attained gaps.

\begin{longtable}[]{@{}
  >{\raggedright\arraybackslash}p{(\linewidth - 4\tabcolsep) * \real{0.3333}}
  >{\raggedleft\arraybackslash}p{(\linewidth - 4\tabcolsep) * \real{0.3333}}
  >{\raggedleft\arraybackslash}p{(\linewidth - 4\tabcolsep) * \real{0.3333}}@{}}
\caption{Observed application time and sampled process-tree
memory.}\label{tbl-pyomo-outcomes}\tabularnewline
\toprule\noalign{}
\begin{minipage}[b]{\linewidth}\raggedright
Implementation
\end{minipage} & \begin{minipage}[b]{\linewidth}\raggedleft
Application end-to-end, min
\end{minipage} & \begin{minipage}[b]{\linewidth}\raggedleft
Sampled worker-tree peak, GiB
\end{minipage} \\
\midrule\noalign{}
\endfirsthead
\toprule\noalign{}
\begin{minipage}[b]{\linewidth}\raggedright
Implementation
\end{minipage} & \begin{minipage}[b]{\linewidth}\raggedleft
Application end-to-end, min
\end{minipage} & \begin{minipage}[b]{\linewidth}\raggedleft
Sampled worker-tree peak, GiB
\end{minipage} \\
\midrule\noalign{}
\endhead
\bottomrule\noalign{}
\endlastfoot
Pyomo/APPSI Gurobi & 55.621 & 55.390 \\
Pyomo/APPSI HiGHS & 441.202 & 55.218 \\
\end{longtable}

Sampling can miss an instantaneous peak; these values are distinct from
requested memory and native solver-managed memory.

Gurobi's lower bound was 12.974647770397 trillion BRL; the HiGHS lower
bound was 12.974647770397 trillion BRL at the shown precision. HiGHS's
smaller incumbent is compatible with the Gurobi interval. Its much
tighter attained gap does not imply that Gurobi would need its reported
runtime to attain that same quality, or that HiGHS would need its full
observed time to attain only the 1\% threshold. The comparison reports
the outcomes actually retained, rather than interpolating unobserved
stopping times.

The HiGHS economic subtotal is 393.331 billion BRL, domestic service is
92.3284\%, and artificial penalties account for 96.9685\% of its
objective. Thus, a near-zero scalar gap coexists with substantial
modeled shortages and capacity exceedances. Objective precision,
financial interpretation and physical service answer different
questions.

The HiGHS log records one processed node and 2,901,682 LP iterations.
Its four-thread ceiling coexisted with one maximum search worker and
parallel search off. These observations illustrate why a thread setting
is not itself a measure of parallel work. The native Gurobi cohort in
the preceding section uses a different implementation and original time
profile; it is not pooled with the repeated common-Pyomo campaign.

The
\href{https://github.com/unb-lamfo-or-ai-hpc/model-agrologistic/blob/ad2ded8132bf5171be65c0113b3d0026c22b2e7b/docs/evidence/scalar_highs_h215_2226536_inspection.json}{retained
HiGHS input and validation record} and the original source-bound
experiment records retain the certificate scope. Appendix E presents the
corresponding accepted investment and flow summaries. The new
common-Pyomo scaling cohort is reported separately below.

\subsection{Computational growth across warehouse
populations}\label{computational-growth-across-warehouse-populations}

The native Gurobi records separate two changes in the workload. At h215
W, the move from DET to STO expands the matrix from 1.56 to 14.05
million variables and from 5.79 to 52.07 million nonzeros, about
ninefold in both cases. The 307 binary variables are shared investment
decisions in STO. With four threads in both observations, native time
rises from 2.04 to 40.76 min and sampled application RSS from 5.54 to
43.13 GiB. The attained gaps are 0.000812\% and 0.740159\%,
respectively. These are observed effort ratios for different
formulations and attained gaps.

Across the W STO series, h215 to h500 increases the number of variables
by 4.51 times and matrix nonzeros by 4.69 times. Model build time grows
by 4.56 times, native time by 3.13 times and sampled application RSS by
2.59 times. The h215 observation used four threads. The h300, h400 and
h500 observations used one. Over this one-thread subset, h300 to h500
increases variables by 2.53 times, matrix nonzeros by 2.58 times, native
time by 2.71 times and RSS by 2.41 times. The h500 incumbent satisfies
the target with an attained gap of 0.207432\% after 127.58 min of native
time and 152.72 min end-to-end. The scale-up comparisons use the
historical native cohort in Table~\ref{tbl-scalar-size},
Table~\ref{tbl-scalar-resources} and Table~\ref{tbl-scalar-quality}.

\subsection{Fixed-instance scaling and attainment in common
Pyomo}\label{fixed-instance-scaling-and-attainment-in-common-pyomo}

All 30 h300 W STO cells returned terminal observations. The two solvers
received the same ordered Pyomo matrix, with 25.11 million variables and
94.97 million nonzeros. The comparison uses three temporal repetitions
per solver and thread ceiling, the common 1\% relative-gap target and an
eight-hour native time limit. Table~\ref{tbl-h300-repeated} summarizes
the median terminal clocks and the count of independently validated
incumbents meeting the target. Appendix C reports the full ranges, gaps
and sampled worker-tree memory.

\begin{longtable}[]{@{}
  >{\raggedleft\arraybackslash}p{(\linewidth - 8\tabcolsep) * \real{0.2000}}
  >{\raggedleft\arraybackslash}p{(\linewidth - 8\tabcolsep) * \real{0.2000}}
  >{\raggedleft\arraybackslash}p{(\linewidth - 8\tabcolsep) * \real{0.2000}}
  >{\raggedleft\arraybackslash}p{(\linewidth - 8\tabcolsep) * \real{0.2000}}
  >{\raggedleft\arraybackslash}p{(\linewidth - 8\tabcolsep) * \real{0.2000}}@{}}
\caption{Common-Pyomo h300 W STO attainment and median times to
termination. HiGHS medians at 2--16 threads describe time-limit returns
without a usable incumbent.}\label{tbl-h300-repeated}\tabularnewline
\toprule\noalign{}
\begin{minipage}[b]{\linewidth}\raggedleft
Threads
\end{minipage} & \begin{minipage}[b]{\linewidth}\raggedleft
Gurobi at target
\end{minipage} & \begin{minipage}[b]{\linewidth}\raggedleft
Gurobi native / end-to-end, min
\end{minipage} & \begin{minipage}[b]{\linewidth}\raggedleft
HiGHS at target
\end{minipage} & \begin{minipage}[b]{\linewidth}\raggedleft
HiGHS native / end-to-end, min
\end{minipage} \\
\midrule\noalign{}
\endfirsthead
\toprule\noalign{}
\begin{minipage}[b]{\linewidth}\raggedleft
Threads
\end{minipage} & \begin{minipage}[b]{\linewidth}\raggedleft
Gurobi at target
\end{minipage} & \begin{minipage}[b]{\linewidth}\raggedleft
Gurobi native / end-to-end, min
\end{minipage} & \begin{minipage}[b]{\linewidth}\raggedleft
HiGHS at target
\end{minipage} & \begin{minipage}[b]{\linewidth}\raggedleft
HiGHS native / end-to-end, min
\end{minipage} \\
\midrule\noalign{}
\endhead
\bottomrule\noalign{}
\endlastfoot
1 & 3/3 & 48.23 / 91.97 & 3/3 & 401.43 / 443.35 \\
2 & 3/3 & 66.18 / 108.94 & 0/3 & 480.79 / 513.14 \\
4 & 3/3 & 65.14 / 113.14 & 0/3 & 480.61 / 513.04 \\
8 & 3/3 & 93.29 / 136.00 & 0/3 & 480.62 / 513.16 \\
16 & 3/3 & 96.39 / 141.15 & 0/3 & 480.60 / 512.78 \\
\end{longtable}

At one thread, the HiGHS-to-Gurobi ratio of median terminal time is 8.32
for the native clock and 4.82 end-to-end. The respective attained gaps
are 0.006879\% and 0.004105\%. At the other four settings, all 12 HiGHS
repetitions reached the limit without a usable incumbent. For these
cells the time to a validated target is unobserved and the
incumbent-based gaps are unavailable. Native times slightly above 480
min are observed terminal readings under the configured budget.

All Gurobi runs reached the target, with attained gaps between
0.004105\% and 0.004122\%. Table~\ref{tbl-h300-scaling} gives the
strong-scaling indicators for the fixed problem under this stopping
policy. They use the median of three runs per ceiling and the
corresponding one-thread median. Efficiency is normalized by the solver
thread ceiling. The observed clocks measure termination; the time of
first reaching the target was not recorded.

\begin{longtable}[]{@{}
  >{\raggedleft\arraybackslash}p{(\linewidth - 8\tabcolsep) * \real{0.2000}}
  >{\raggedleft\arraybackslash}p{(\linewidth - 8\tabcolsep) * \real{0.2000}}
  >{\raggedleft\arraybackslash}p{(\linewidth - 8\tabcolsep) * \real{0.2000}}
  >{\raggedleft\arraybackslash}p{(\linewidth - 8\tabcolsep) * \real{0.2000}}
  >{\raggedleft\arraybackslash}p{(\linewidth - 8\tabcolsep) * \real{0.2000}}@{}}
\caption{Descriptive Gurobi h300 strong-scaling indicators based on
median terminal times under the shared stopping
policy.}\label{tbl-h300-scaling}\tabularnewline
\toprule\noalign{}
\begin{minipage}[b]{\linewidth}\raggedleft
Threads
\end{minipage} & \begin{minipage}[b]{\linewidth}\raggedleft
Native speedup
\end{minipage} & \begin{minipage}[b]{\linewidth}\raggedleft
Native efficiency, \%
\end{minipage} & \begin{minipage}[b]{\linewidth}\raggedleft
End-to-end speedup
\end{minipage} & \begin{minipage}[b]{\linewidth}\raggedleft
End-to-end efficiency, \%
\end{minipage} \\
\midrule\noalign{}
\endfirsthead
\toprule\noalign{}
\begin{minipage}[b]{\linewidth}\raggedleft
Threads
\end{minipage} & \begin{minipage}[b]{\linewidth}\raggedleft
Native speedup
\end{minipage} & \begin{minipage}[b]{\linewidth}\raggedleft
Native efficiency, \%
\end{minipage} & \begin{minipage}[b]{\linewidth}\raggedleft
End-to-end speedup
\end{minipage} & \begin{minipage}[b]{\linewidth}\raggedleft
End-to-end efficiency, \%
\end{minipage} \\
\midrule\noalign{}
\endhead
\bottomrule\noalign{}
\endlastfoot
1 & 1.000 & 100.0 & 1.000 & 100.0 \\
2 & 0.729 & 36.4 & 0.844 & 42.2 \\
4 & 0.740 & 18.5 & 0.813 & 20.3 \\
8 & 0.517 & 6.5 & 0.676 & 8.5 \\
16 & 0.500 & 3.1 & 0.652 & 4.1 \\
\end{longtable}

Each multithreaded Gurobi setting takes longer than the one-thread
setting on both median clocks. At 16 threads the median native time is
approximately double and the end-to-end time is 1.53 times the
one-thread value. The sampled worker-tree peak rises from 74.42 GiB at
one thread to 100.88 GiB at 16. At two and four threads the native-time
ranges overlap the one-thread range. At eight and sixteen threads they
lie entirely above it in the three observed repetitions. The per-setting
ranges characterize this run block.

\subsection{Evidence at h400 and coverage of
h500}\label{evidence-at-h400-and-coverage-of-h500}

The h400 common-Pyomo matrix contains 42.11 million variables and 161.25
million nonzeros, compared with 25.11 million and 94.97 million at h300.
The E1 and E2 follow-up runs complete one observation per solver and
thread ceiling at h400 W STO. The ten receipts bind the same inputs,
ordered matrix, bounds and hardware class to the 192 GiB profile and
eight-hour native budget. E1 supplies one and sixteen threads, and E2
supplies two, four and eight. Appendix C distinguishes these
observations from the original h400 pilots.

Gurobi reached the 1\% target in all five configurations. Each returned
the same reported penalized objective, BRL 7.593395 trillion, and
economic subtotal, BRL 381.573 billion. The attained gaps agree to six
decimal places at 0.836784\%. HiGHS reached the time limit without a
usable incumbent at every thread ceiling. Its native terminal times
range from 481.13 to 481.39 min. Its incumbent objectives and gaps are
unavailable. Table~\ref{tbl-h400-followup} reports both terminal clocks
and Gurobi's descriptive speedup relative to the E1 one-thread
observation.

\begin{longtable}[]{@{}
  >{\raggedleft\arraybackslash}p{(\linewidth - 6\tabcolsep) * \real{0.2500}}
  >{\raggedleft\arraybackslash}p{(\linewidth - 6\tabcolsep) * \real{0.2500}}
  >{\raggedleft\arraybackslash}p{(\linewidth - 6\tabcolsep) * \real{0.2500}}
  >{\raggedleft\arraybackslash}p{(\linewidth - 6\tabcolsep) * \real{0.2500}}@{}}
\caption{h400 W STO terminal clocks and Gurobi speedup in the E1--E2
follow-up. All Gurobi runs reached the target. All HiGHS times describe
censoring without a usable incumbent. Each setting has one
observation.}\label{tbl-h400-followup}\tabularnewline
\toprule\noalign{}
\begin{minipage}[b]{\linewidth}\raggedleft
Threads
\end{minipage} & \begin{minipage}[b]{\linewidth}\raggedleft
Gurobi native / end-to-end, min
\end{minipage} & \begin{minipage}[b]{\linewidth}\raggedleft
Gurobi native / end-to-end speedup
\end{minipage} & \begin{minipage}[b]{\linewidth}\raggedleft
HiGHS native / end-to-end, min
\end{minipage} \\
\midrule\noalign{}
\endfirsthead
\toprule\noalign{}
\begin{minipage}[b]{\linewidth}\raggedleft
Threads
\end{minipage} & \begin{minipage}[b]{\linewidth}\raggedleft
Gurobi native / end-to-end, min
\end{minipage} & \begin{minipage}[b]{\linewidth}\raggedleft
Gurobi native / end-to-end speedup
\end{minipage} & \begin{minipage}[b]{\linewidth}\raggedleft
HiGHS native / end-to-end, min
\end{minipage} \\
\midrule\noalign{}
\endhead
\bottomrule\noalign{}
\endlastfoot
1 & 87.36 / 162.60 & 1.000 / 1.000 & 481.13 / 534.88 \\
2 & 80.03 / 154.83 & 1.092 / 1.050 & 481.39 / 535.59 \\
4 & 139.63 / 214.98 & 0.626 / 0.756 & 481.35 / 536.08 \\
8 & 150.36 / 226.35 & 0.581 / 0.718 & 481.37 / 536.96 \\
16 & 134.60 / 207.78 & 0.649 / 0.783 & 481.28 / 534.81 \\
\end{longtable}

The two-thread Gurobi run had the lowest observed time on both clocks,
with nominal efficiencies of 54.6\% native and 52.5\% end-to-end. Its
sampled worker-tree peak was 133.13 GiB compared with 113.31 GiB at one
thread. Across both solvers the largest sampled worker-tree and
job-cgroup peaks were 162.99 and 163.45 GiB, respectively, both from
Gurobi at sixteen threads. Both were within the follow-up gate limits.
Appendix C reports the two memory measurements for each cell.

The original 90-cell campaign remains a separate record. It retained 20
validated target attainments and 14 observations without usable
incumbents. The remaining 56 cells, including 26 at h400 and all 30 at
h500, were cancelled before execution. The native h500 solution remains
the largest validated instance in the reported series.

\subsection{Planning diagnostics from the first h300
pair}\label{planning-diagnostics-from-the-first-h300-pair}

The two returned h300 W STO pilots use Gurobi 13.0.3 and HiGHS 1.15.1
through Pyomo/APPSI. Both original incumbents passed independent primal
and cost validation and the prescribed 1\% relative-gap criterion. They
share the same frozen workbook, ordered warehouse selection, selected
directed graph, scenario probabilities, mathematical contract,
bounded-slack policy and ordered Pyomo algebra. The hardware-class
fingerprints and requested resource profiles also agree. This
establishes a matched input and formulation basis; it does not require
identical native presolve or search trajectories.

For the first common-Pyomo h300 pair,
Table~\ref{tbl-scaling-h300-settings} identifies the settings,
Table~\ref{tbl-scaling-h300-pilot} and
Table~\ref{tbl-scaling-h300-resources} summarize computational outcomes,
and Table~\ref{tbl-scaling-h300-physical} presents operating
diagnostics.

\begin{longtable}[]{@{}
  >{\raggedright\arraybackslash}p{(\linewidth - 4\tabcolsep) * \real{0.3333}}
  >{\raggedright\arraybackslash}p{(\linewidth - 4\tabcolsep) * \real{0.3333}}
  >{\raggedright\arraybackslash}p{(\linewidth - 4\tabcolsep) * \real{0.3333}}@{}}
\caption{Settings of the first matched h300
pair.}\label{tbl-scaling-h300-settings}\tabularnewline
\toprule\noalign{}
\begin{minipage}[b]{\linewidth}\raggedright
Experimental control
\end{minipage} & \begin{minipage}[b]{\linewidth}\raggedright
Gurobi
\end{minipage} & \begin{minipage}[b]{\linewidth}\raggedright
HiGHS
\end{minipage} \\
\midrule\noalign{}
\endfirsthead
\toprule\noalign{}
\begin{minipage}[b]{\linewidth}\raggedright
Experimental control
\end{minipage} & \begin{minipage}[b]{\linewidth}\raggedright
Gurobi
\end{minipage} & \begin{minipage}[b]{\linewidth}\raggedright
HiGHS
\end{minipage} \\
\midrule\noalign{}
\endhead
\bottomrule\noalign{}
\endlastfoot
Interface & Pyomo/APPSI & Pyomo/APPSI \\
Native time allowance, s & 28,800 & 28,800 \\
Relative-gap target & 0.01 & 0.01 \\
Solver threads & 1 & 1 \\
Seed & Seed=42 & random\_seed=42 \\
Numerical controls & FeasibilityTol=IntFeasTol=1e-9 &
mip\_feasibility\_tolerance=1e-9 \\
Search control & Automatic methods/presolve/cuts; Heuristics=0.05 &
parallel=on \\
\end{longtable}

With one permitted thread, enabling HiGHS parallel search does not
establish multiworker parallel execution.

\begin{longtable}[]{@{}
  >{\raggedright\arraybackslash}p{(\linewidth - 6\tabcolsep) * \real{0.2500}}
  >{\raggedleft\arraybackslash}p{(\linewidth - 6\tabcolsep) * \real{0.2500}}
  >{\raggedleft\arraybackslash}p{(\linewidth - 6\tabcolsep) * \real{0.2500}}
  >{\raggedleft\arraybackslash}p{(\linewidth - 6\tabcolsep) * \real{0.2500}}@{}}
\caption{Computational outcomes for the first h300 observations under
common-Pyomo protocol.}\label{tbl-scaling-h300-pilot}\tabularnewline
\toprule\noalign{}
\begin{minipage}[b]{\linewidth}\raggedright
Solver / threads / repetition
\end{minipage} & \begin{minipage}[b]{\linewidth}\raggedleft
Objective, BRL trillion
\end{minipage} & \begin{minipage}[b]{\linewidth}\raggedleft
Common gap, \%
\end{minipage} & \begin{minipage}[b]{\linewidth}\raggedleft
Native runtime, min
\end{minipage} \\
\midrule\noalign{}
\endfirsthead
\toprule\noalign{}
\begin{minipage}[b]{\linewidth}\raggedright
Solver / threads / repetition
\end{minipage} & \begin{minipage}[b]{\linewidth}\raggedleft
Objective, BRL trillion
\end{minipage} & \begin{minipage}[b]{\linewidth}\raggedleft
Common gap, \%
\end{minipage} & \begin{minipage}[b]{\linewidth}\raggedleft
Native runtime, min
\end{minipage} \\
\midrule\noalign{}
\endhead
\bottomrule\noalign{}
\endlastfoot
Gurobi / 1 / 1 & 9.289109 & 0.004105 & 48.229 \\
HiGHS / 1 / 1 & 9.289359 & 0.006879 & 401.435 \\
\end{longtable}

Attained quality and native time for the first h300 observations under
the new common-Pyomo protocol. Both original validations were accepted.

\begin{longtable}[]{@{}lrr@{}}
\caption{Application effort and observed memory scope for the same pair.
}\label{tbl-scaling-h300-resources}\tabularnewline
\toprule\noalign{}
Solver & Application end-to-end, min & Sampled worker-tree peak, GiB \\
\midrule\noalign{}
\endfirsthead
\toprule\noalign{}
Solver & Application end-to-end, min & Sampled worker-tree peak, GiB \\
\midrule\noalign{}
\endhead
\bottomrule\noalign{}
\endlastfoot
Gurobi & 91.794 & 74.535 \\
HiGHS & 443.350 & 84.820 \\
\end{longtable}

The lower bounds are 9.288727295144 and 9.288719505216 trillion BRL for
Gurobi and HiGHS, respectively. Their objective-bound intervals overlap.
Gurobi returned a smaller incumbent and a tighter attained gap in this
pair. Both gaps are substantially below the common 1\% stopping target,
but neither is an exact-zero certificate. The observed native-time
ratio, HiGHS/Gurobi, is 8.32. This is a descriptive ratio for two
retained outcomes, not a causal speedup or an estimate of runtime
variability from repetitions.

The shared ordered algebra has 25,107,544 variables, including 392
binaries, 1,056,962 linear constraints and 94,967,766 nonzeros. Each
solver processed one node; Gurobi and HiGHS report 3,355,126 and
3,998,477 LP iterations. Construction, solver transfer, matrix
identification and extraction add effort outside native optimization.
Some clocks enclose others and must not be added as disjoint components.
Sampled process-tree memory, Slurm MaxRSS and the 192 GiB requested
allowance retain their separate meanings; the observations do not
certify production isolation or an unobserved instantaneous peak.

\begin{longtable}[]{@{}lrr@{}}
\caption{Physical and monetary diagnostics from the original accepted
exports. }\label{tbl-scaling-h300-physical}\tabularnewline
\toprule\noalign{}
Planning diagnostic & Gurobi & HiGHS \\
\midrule\noalign{}
\endfirsthead
\toprule\noalign{}
Planning diagnostic & Gurobi & HiGHS \\
\midrule\noalign{}
\endhead
\bottomrule\noalign{}
\endlastfoot
Financial subtotal, BRL billion & 385.502 & 385.486 \\
Artificial penalties, \% of objective & 95.850 & 95.850 \\
Expected domestic service, \% & 91.3256 & 91.3252 \\
Unmet demand, million t & 58.787 & 58.790 \\
Static exceedance, million tonne-periods & 46.444 & 46.444 \\
Reception overflow, million t & 33.839 & 33.840 \\
\end{longtable}

HiGHS has a slightly lower financial subtotal while Gurobi has the lower
penalized total. The financial subtotal alone therefore does not order
these solutions under the modeled objective. Likewise, near agreement in
aggregate service and costs does not imply identical investments or
flows. Appendix E retains their separate first-stage lists and directed
flow-distance companions. Stock exposure and reception overflow have
different dimensions and are not added as installed emergency capacity.

The
\href{https://github.com/unb-lamfo-or-ai-hpc/model-agrologistic/blob/083ef96cc99a346204df22551ce0739c10e8a2f7/docs/evidence/paired_h300_first_pair_20261101.json}{source-bound
paired audit} retains original archive identities, matching
input/algebra records, effective parameters, bounds and measurement
scopes. Appendix C identifies this pair within the h300 repetition grid
and gives the full timing, gap and memory ranges for each setting.


\subsection{Controlled h300 root-relaxation results}\label{exploratory-h300-root-relaxation-results}

The controlled h300 STO/W assessment varied four root-relaxation methods, three heuristic profiles and five thread ceilings on distinct physical cores, with an eight-hour native time budget and 192 GiB per job. The first 60-configuration block was closed with 33 independently validated incumbents, nine time-limited returns without a usable incumbent, seven native memory-limit returns and seven job-memory-guard stops. One execution was cancelled during optimization and three configurations were cancelled before starting. The two additional temporal blocks were not executed. Table~\ref{tbl-root-partial} summarizes the resulting coverage; Appendix C11 gives individual configurations, original validation outcomes and resource measurements.

\begin{table}[htbp]
\centering\small
\caption{Outcomes of the closed first h300 block. Native memory denotes Gurobi's soft memory limit; guard denotes the external job-memory monitor. Cancellation includes one interrupted and three unstarted configurations. Counts describe this block, without estimating success probabilities.}
\label{tbl-root-partial}
\begin{tabular}{lrrrrr}
\toprule
Root LP method & Accepted & Time & Native memory & Guard & Cancelled\\
\midrule
Dual simplex (1) & 10 & 0 & 2 & 1 & 2\\
Barrier (2) & 10 & 0 & 2 & 2 & 1\\
Deterministic concurrent (4) & 11 & 0 & 2 & 2 & 0\\
PDHG (6) & 2 & 9 & 1 & 2 & 1\\
\midrule
Total & 33 & 9 & 7 & 7 & 4\\
\bottomrule
\end{tabular}
\end{table}

Figure~\ref{fig-root-final} separates the completed root-relaxation clock from the native terminal clock for accepted observations under the standard heuristic profile. Barrier root time decreased from 122.3 minutes at one core to 61.3 minutes at sixteen cores; native terminal time decreased from 137.8 to 78.1 minutes. These observations correspond to a root-phase time ratio of 1.99 and an observed terminal speedup of 1.77, with nominal terminal efficiency of 11.0\% at sixteen cores. The four-core barrier configuration was cancelled before execution. Each available configuration has one observation, so these values describe the observed fixed-instance response under the stated protocol.

\begin{figure}[htbp]
\centering
\includegraphics[width=\linewidth]{h300_root_final.pdf}
\caption{Completed root-relaxation and native terminal times for accepted h300 STO/W observations with standard heuristics. Methods 1, 2, 4 and 6 denote dual simplex, barrier, deterministic concurrent optimization and PDHG. Each point represents one configuration; line segments connect adjacent available core settings. Unshown settings have no accepted result in this block; Appendix C11 retains their outcomes.}
\label{fig-root-final}
\end{figure}

Dual simplex returned accepted solutions in 41.1 minutes at one core and 50.1 minutes at sixteen cores. Deterministic concurrent optimization took 41.7 and 66.5 minutes at these settings. Thus, the method with the clearest observed response to additional cores was slower in absolute time than the accepted one-core dual-simplex solution. PDHG returned accepted standard-profile solutions at eight and sixteen cores in 360.4 and 312.9 minutes; its standard-profile runs at one, two and four cores reached the time limit without a usable incumbent.

Complete logs resolve the seven returns labelled as errors by the application: Gurobi reported \emph{Memory limit reached} with no incumbent. All occurred with the general heuristics disabled at eight or sixteen cores. The seven external memory-guard stops occurred with the NoRel profile at eight or sixteen cores. Their retained cgroup counters show zero OOM-kill and failure-count increments. Native memory accounting, job-cgroup memory and the Slurm allocation have distinct scopes. Slurm reported 49 completed tasks, including the nine time-limited and seven native memory-limited returns, whereas 33 had accepted solutions.

Across those 33 solutions, relative gaps ranged from 0.001350\% to 0.006175\%, and the total penalized objective ranged from BRL 9.288940 to 9.289372 trillion. Economic cost ranged from BRL 385.343 to 385.502 billion, while domestic service remained between 91.3237\% and 91.3258\%. Artificial penalties accounted for approximately 95.85\% of the objective. The different solution methods therefore changed computational effort and attainment more substantially than these aggregate agricultural outcomes. Accepted runs explored between one and 47 nodes, showing that the earlier one-node behavior depends on the solver configuration and trajectory.

\section{Discussion}\label{discussion}

\subsection{What scales with the planning
problem}\label{what-scales-with-the-planning-problem}

The h215 results show the cost of introducing scenario-dependent
operation while retaining the same shared investment binaries. The
near-ninefold expansion in variables and nonzeros accompanies a
twentyfold increase in the observed native time. Matrix size gives a
direct measure of formulation size and helps explain storage and build
requirements. Search and LP work also depend on the coefficients,
bounds, incumbent path and attained stopping gap.

In the native one-thread STO series, the increasing build times and RSS
from h300 to h500 are directly relevant to sizing a workstation or HPC
allocation. The higher effort buys a different planning choice set, with
additional candidate sites and different admissible routes. Population,
arcs and available decisions change jointly. This is problem-size
scale-up. Classical weak scaling would require a design that maintains
approximately constant workload per resource.

\subsection{Thread ceilings, terminal effort and usable
plans}\label{thread-ceilings-terminal-effort-and-usable-plans}

The common-Pyomo h300 block shows why a solver thread setting should be
evaluated jointly with the quality of the returned solution. The Gurobi
comparison has complete target attainment and closely aligned attained
gaps across the five settings. One thread delivered the lowest median
time on both clocks, with less sampled memory than the higher settings.
The gap between native and end-to-end ratios matters for practitioners
who need a validated investment report. Build, transfer and checking
contribute to that elapsed time. The observations measure neither CPU
utilization nor a causal allocation of time to parallel overhead.

At one thread both solvers delivered usable h300 plans. At the higher
ceilings, HiGHS returned only a bound within the observed budget.
Treating those terminal times as times to solution would reverse the
scientific meaning of the comparison. The useful result is the observed
attainment contrast at a fixed size, together with the bound,
termination state and resource use. These fixed-seed repetitions
characterize variability on the same instance and protocol. They do not
sample independent planning problems or establish a general solver
ranking.

At h400, the two-thread Gurobi observation took 8.4\% less native time
and 4.8\% less end-to-end time than the one-thread observation, while
the sampled worker peak increased by 19.82 GiB. The 4-, 8- and 16-thread
observations were slower on both clocks. The repeated h300 grid and the
single-repetition h400 grid therefore give different rankings of the
one- and two-thread settings. Repeated h400 runs would be needed to
estimate whether its observed two-thread advantage persists. HiGHS
reached the budget without a usable incumbent at all five ceilings, so
the h400 solver contrast concerns attainment as well as effort.

The original h400 gate rejected progression because the sampled
worker-tree peak exceeded its ceiling, although the incumbent met the
numerical target. The subsequent E1--E2 runs passed their own worker and
cgroup criteria at the same 192 GiB allocation. The largest sampled
worker peak was within 0.21 GiB of its follow-up ceiling. Appendix C
gives the scope and limits for each observation. A completed single
repetition provides evidence for those runs and leaves the memory
requirements of repetitions and h500 open.

The closed root-method block locates an important part of the computational response in the initial LP relaxation. With standard heuristics, barrier reduced root time as physical cores increased, while dual simplex and deterministic concurrent optimization had higher terminal times at sixteen cores than at one. Barrier's greater parallel responsiveness coexisted with a longer absolute solution time than one-core dual simplex. Selecting a solver configuration therefore requires considering elapsed time, memory and accepted solution quality together. The nine time-limited PDHG returns and the two distinct memory protections explain why terminal runtime alone cannot rank all configurations. These observations characterize one fixed instance and one temporal block; the omitted repetitions leave temporal repeatability unmeasured.

\subsection{Implications for grain-storage
investment}\label{implications-for-grain-storage-investment}

The native W STO plans show a shift in expenditure as the candidate set
grows. From h215 to h500, the returned one-time investment increases
from BRL 9.82 billion to 14.32 billion while transport decreases from
BRL 367.14 billion to 346.07 billion. The economic subtotal falls from
BRL 393.90 billion to 373.17 billion. Installed candidate capacity
increases from 7.38 to 10.99 million tonnes. The h500 incumbent
allocates capacity across all 354 candidate locations. That distribution
reflects the encoded costs, routing and service trade-offs and would
require site-level validation before implementation.

The h215 W/D results illustrate why economic expenditure and service
should be read together. With direct deliveries, the returned D plans
have economic subtotals 13.11\% lower in DET and 13.66\% lower in STO.
Domestic service falls by 1.97 and 2.41 percentage points, respectively.
Direct deliveries bypass some warehouse handling and stock, while the
optimization trades off shortage and capacity penalties under the same
weights. In the native W STO series, service also varies
nonmonotonically as the candidate set expands. Lower expenditure alone
is therefore an incomplete basis for selecting an investment program.

Artificial penalties represent roughly 94--97\% of the native W/D
penalized objectives. A certified gap applies to the penalized total.
The economic subtotal, delivered grain and the three physical slack
measures remain separate. A planner can use those diagnostics to
identify where a returned investment plan depends on unmet demand or
capacity exceedance. Stock summed over time is measured in
tonne-periods, whereas reception overflow is measured in tonnes.

The first common-Pyomo h300 pair further shows how similar aggregate
outcomes can hide different investments. Both solutions open the same
154 candidates and report domestic service close to 91.33\%. They share
69 expansion sites and 76 bulk-handling sites. The Gurobi plan adds one
expansion to the shared set. The HiGHS plan adds one bulk-handling
conversion to the shared set. Appendix E retains the individual facility
identifiers, capacities and directed flows, which are the evidence
needed to identify decisions that are consistent between the two plans.

\subsection{Relationship to agricultural computing
research}\label{relationship-to-agricultural-computing-research}

The evidence complements the tactical effort measurements of
\citet{reis2023} by linking construction, solver effort and memory to
long-lived warehouse investment. The discrete sizing and demand
fulfillment in \citet{pourmohammadi2020} frame the different
implications of continuous sizing and penalized domestic shortage here.
The capacity-relief mechanism in \citet{milanez2026} motivates reporting
operational quantities alongside the numerical certificate. The
distributed decomposition in \citet{shastri2011} addresses a related
computational burden using a different solution strategy. The present
measurements characterize direct extensive-form solution under declared
configurations and supply a reproducible comparison within the study.

The complementary human-supervised workflow links each computational
claim to an input, execution and validation record. Appendix B describes
the decisions and checks that make this evidence traceable. The primary
contribution remains the joint characterization of computational
requirements and grain-storage decisions.

\section{Final considerations, limitations and future
work}\label{final-considerations-limitations-and-future-work}

Eleven validated native Gurobi observations met the 1\% gap target. The
native W STO series reached 63.43 million variables and 244.04 million
nonzeros at h500. Its returned plan assigns capacity to 354 candidate
sites, expands 72 existing warehouses and converts 74 for bulk handling.
The attained gap is 0.207432\% and modeled domestic service is 91.636\%.
The reported investment, costs and physical diagnostics make the
planning trade-offs inspectable.

Gurobi and HiGHS returned independently validated Pyomo plans for h215 W
STO and for the first matched h300 W STO pair. Across the complete h300
repetition grid, Gurobi met the target in 15 of 15 runs and HiGHS in
three of 15. The Gurobi one-thread median was lowest on both timing
clocks. At 16 threads the nominal native efficiency was 3.1\% relative
to the one-thread baseline. These indicators describe a fixed-instance
comparison with the declared options and gap policy. The complete h400
E1--E2 grid contains one observation per setting. Gurobi attained the
target in all five runs and HiGHS reached its time limit without a
usable incumbent in all five. The fastest observed Gurobi setting was
two threads. The h500 computational result comes from the historical
native implementation.

The controlled h300 block closed with 33 accepted solutions, nine time-limited returns, fourteen memory-related stops and four cancellations. It provides conditional evidence of barrier's fixed-instance parallel response and of the different behavior of the other root methods. The study contains one temporal block; two planned repetitions were not executed. The observed one-node behavior is not a structural limit on branch-and-bound.

The main limitations are the fixed sets of instances and scenarios,
three temporal repetitions per repeated setting, finite compute budgets
and incomplete coverage at the larger sizes. The investment conclusions
also depend on the represented sites, costs, period structure and
penalties. Decision support for implementation requires credible local
costs, site suitability and explicit service requirements.

\subsection{Future work}\label{future-work}

Future studies should assess the repeatability of the method-specific root-time response, examine formulations or decomposition strategies that reduce the extensive-form memory burden, and evaluate service, export and capacity scenarios against an explicitly frozen effective baseline. These investigations extend the closed MVP2 assessment and require separate experimental protocols. Agricultural improvement should be measured through admissible investment and operating plans, domestic service and physical exceedances alongside economic cost.


Repeat the h400 thread grid under comparable conditions to quantify
variability in time, attainment and memory headroom. Progression to h500
depends on a separate resource qualification. Retain the matrix,
stopping, timing and memory identities needed for each comparison, as
specified in Appendix C.

A separate planning study could test calibrated costs and service
requirements. That comparison would evaluate how the resulting
investments change relative to the present plans.

\section{Reproducibility and
availability}\label{reproducibility-and-availability}

The source code, experiment definitions and numerical records are
available in
\href{https://github.com/unb-lamfo-or-ai-hpc/model-agrologistic}{GitHub
repository}. Each experimental record identifies the source version,
input, warehouse population, settings, decisions and original
validation. A later integrity check does not change the identities of
the executed software or its results.

The article and supplementary material are written in Quarto\footnote{\url{https://quarto.org/}}
and can be rendered without an optimizer or HPC credentials. Appendix A
presents the literature review; Appendix B explains AI-assisted research
and verification; Appendix C describes the Python implementations,
independent validation and GAMS\footnote{\url{https://www.gams.com/products/solvers/}}
translation; Appendix D specifies input data and effective parameters.
GAMS compilation and numerical parity have not been verified. The Python
results therefore do not constitute numerical validation of that
translation.

The supplementary material is accessible as
\href{https://unb-lamfo-or-ai-hpc.github.io/model-agrologistic/supplementary-review.html}{HTML}
and
\href{https://unb-lamfo-or-ai-hpc.github.io/model-agrologistic/supplementary-review.pdf}{PDF}.
Exact source snapshots and data identities are given in the linked
repository.

The editable
\href{https://unb-lamfo-or-ai-hpc.github.io/model-agrologistic/index.tex}{article
TeX} and
\href{https://unb-lamfo-or-ai-hpc.github.io/model-agrologistic/supplementary-review.tex}{appendix
TeX} are generated from the Quarto sources. Their figures, bibliography,
class files and notices accompany the
\href{https://unb-lamfo-or-ai-hpc.github.io/model-agrologistic/editable-sources.zip}{complete
editable project}. The individual TeX downloads are intended for
editing; compilation requires those accompanying dependencies.

\subsection{Data availability}\label{data-availability}

The research dataset is publicly available in Zenodo, version 0.2.0
\citep{celestino2026dataset}, at
\href{https://doi.org/10.5281/zenodo.22751909}{10.5281/zenodo.22751909}.
The
\href{https://zenodo.org/records/22751909/files/agrologistic-v020-inputs.zip?download=1}{processed
input archive} contains the four effective workbooks. Appendix D lists
the exact members and checksums for h215, h300, h400 and h500.
Appendices D and E identify the decision vectors and flow-distance
companions and explain how to access them. OpenStreetMap derivatives
retain ODbL terms, and project-authored code is MIT. File-level access
and redistribution conditions are recorded in the dataset inventories.

\section{Acknowledgements}\label{acknowledgements}

The authors gratefully acknowledge the Núcleo de Processamento de Alto
Desempenho (NPAD) at the Universidade Federal do Rio Grande do Norte
(UFRN), Brazil, for providing access to the high-performance computing
infrastructure used in this research.

\section{Declaration of generative AI and AI-assisted technologies in
the manuscript preparation
process}\label{declaration-of-generative-ai-and-ai-assisted-technologies-in-the-manuscript-preparation-process}

During preparation, the authors used Gurobot, Gemini and GPT/ChatGPT
systems, including ChatGPT Codex, for technical discussion, software
development, language refinement and visualization. Historical records
name Gemini 3.1 Pro and ChatGPT Codex with models 6.1 Sol and 6.0 Astra
where documented; this does not attribute all work to a verified
model/version chronology. The methodological section and Appendix B
distinguish retained evidence from retrospective accounts. The human
authors reviewed and edited the content and retain full responsibility
for scientific claims and the article.

\section{References}\label{references}

\renewcommand{\bibsection}{}
\bibliography{references.bib}





\end{document}
